%**************************************%
%*    Generated from PreTeXt source   *%
%*    on 2019-05-29T22:02:12-06:00    *%
%*                                    *%
%*      https://pretextbook.org       *%
%*                                    *%
%**************************************%
\documentclass[oneside,10pt,]{book}
%% Custom Preamble Entries, early (use latex.preamble.early)
%% Default LaTeX packages
%%   1.  always employed (or nearly so) for some purpose, or
%%   2.  a stylewriter may assume their presence
\usepackage{geometry}
%% Some aspects of the preamble are conditional,
%% the LaTeX engine is one such determinant
\usepackage{ifthen}
%% etoolbox has a variety of modern conveniences
\usepackage{etoolbox}
\usepackage{ifxetex,ifluatex}
%% Raster graphics inclusion
\usepackage{graphicx}
%% Color support, xcolor package
%% Always loaded, for: add/delete text, author tools
%% Here, since tcolorbox loads tikz, and tikz loads xcolor
\PassOptionsToPackage{usenames,dvipsnames,svgnames,table}{xcolor}
\usepackage{xcolor}
%% Colored boxes, and much more, though mostly styling
%% skins library provides "enhanced" skin, employing tikzpicture
%% boxes may be configured as "breakable" or "unbreakable"
%% "raster" controls grids of boxes, aka side-by-side
\usepackage{tcolorbox}
\tcbuselibrary{skins}
\tcbuselibrary{breakable}
\tcbuselibrary{raster}
%% We load some "stock" tcolorbox styles that we use a lot
%% Placement here is provisional, there will be some color work also
%% First, black on white, no border, transparent, but no assumption about titles
\tcbset{ bwminimalstyle/.style={size=minimal, boxrule=-0.3pt, frame empty,
colback=white, colbacktitle=white, coltitle=black, opacityfill=0.0} }
%% Second, bold title, run-in to text/paragraph/heading
%% Space afterwards will be controlled by environment,
%% dependent of constructions of the tcb title
\tcbset{ runintitlestyle/.style={fonttitle=\normalfont\bfseries, attach title to upper} }
%% Spacing prior to each exercise, anywhere
\tcbset{ exercisespacingstyle/.style={before skip={1.5ex plus 0.5ex}} }
%% Spacing prior to each block
\tcbset{ blockspacingstyle/.style={before skip={2.0ex plus 0.5ex}} }
%% xparse allows the construction of more robust commands,
%% this is a necessity for isolating styling and behavior
%% The tcolorbox library of the same name loads the base library
\tcbuselibrary{xparse}
%% Hyperref should be here, but likes to be loaded late
%%
%% Inline math delimiters, \(, \), need to be robust
%% 2016-01-31:  latexrelease.sty  supersedes  fixltx2e.sty
%% If  latexrelease.sty  exists, bugfix is in kernel
%% If not, bugfix is in  fixltx2e.sty
%% See:  https://tug.org/TUGboat/tb36-3/tb114ltnews22.pdf
%% and read "Fewer fragile commands" in distribution's  latexchanges.pdf
\IfFileExists{latexrelease.sty}{}{\usepackage{fixltx2e}}
%% Footnote counters and part/chapter counters are manipulated
%% April 2018:  chngcntr  commands now integrated into the kernel,
%% but circa 2018/2019 the package would still try to redefine them,
%% so we need to do the work of loading conditionally for old kernels.
%% From version 1.1a,  chngcntr  should detect defintions made by LaTeX kernel.
\ifdefined\counterwithin
\else
    \usepackage{chngcntr}
\fi
%% Text height identically 9 inches, text width varies on point size
%% See Bringhurst 2.1.1 on measure for recommendations
%% 75 characters per line (count spaces, punctuation) is target
%% which is the upper limit of Bringhurst's recommendations
\geometry{letterpaper,total={340pt,9.0in}}
%% Custom Page Layout Adjustments (use latex.geometry)
%% This LaTeX file may be compiled with pdflatex, xelatex, or lualatex executables
%% LuaTeX is not explicitly supported, but we do accept additions from knowledgeable users
%% The conditional below provides  pdflatex  specific configuration last
%% The following provides engine-specific capabilities
%% Generally, xelatex is necessary non-Western fonts
\ifthenelse{\boolean{xetex} \or \boolean{luatex}}{%
%% begin: xelatex and lualatex-specific configuration
\ifxetex\usepackage{xltxtra}\fi
%% realscripts is the only part of xltxtra relevant to lualatex 
\ifluatex\usepackage{realscripts}\fi
%% fontspec package provides extensive control of system fonts,
%% meaning *.otf (OpenType), and apparently *.ttf (TrueType)
%% that live *outside* your TeX/MF tree, and are controlled by your *system*
%% fontspec will make Latin Modern (lmodern) the default font
\usepackage{fontspec}
%% 
%% Extensive support for other languages
\usepackage{polyglossia}
%% Set main/default language based on pretext/@xml:lang value
%% document language code is "en-US", US English
%% usmax variant has extra hypenation
\setmainlanguage[variant=usmax]{english}
%% Enable secondary languages based on discovery of @xml:lang values
%% Enable fonts/scripts based on discovery of @xml:lang values
%% Western languages should be ably covered by Latin Modern Roman
%% end: xelatex and lualatex-specific configuration
}{%
%% begin: pdflatex-specific configuration
\usepackage[utf8]{inputenc}
%% PreTeXt will create a UTF-8 encoded file
%% begin: font setup and configuration for use with pdflatex
\usepackage{lmodern}
\usepackage[T1]{fontenc}
%% end: font setup and configuration for use with pdflatex
%% end: pdflatex-specific configuration
}
%% Monospace font: Inconsolata (zi4)
%% Sponsored by TUG: http://levien.com/type/myfonts/inconsolata.html
%% See package documentation for excellent instructions
%% One caveat, seem to need full file name to locate OTF files
%% Loads the "upquote" package as needed, so we don't have to
%% Upright quotes might come from the  textcomp  package, which we also use
%% We employ the shapely \ell to match Google Font version
%% pdflatex: "varqu" option produces best upright quotes
%% xelatex,lualatex: add StylisticSet 1 for shapely \ell
%% xelatex,lualatex: add StylisticSet 2 for plain zero
%% xelatex,lualatex: we add StylisticSet 3 for upright quotes
%% 
\ifthenelse{\boolean{xetex} \or \boolean{luatex}}{%
%% begin: xelatex and lualatex-specific monospace font
\usepackage{zi4}
\setmonofont[BoldFont=Inconsolatazi4-Bold.otf,StylisticSet={1,3}]{Inconsolatazi4-Regular.otf}
%% end: xelatex and lualatex-specific monospace font
}{%
%% begin: pdflatex-specific monospace font
%% "varqu" option provides textcomp \textquotedbl glyph
%% "varl"  option provides shapely "ell"
\usepackage[varqu,varl]{zi4}
%% end: pdflatex-specific monospace font
}
%% \mono macro for content of "c" element, and XML parts
\newcommand{\mono}[1]{\texttt{#1}}
%% Symbols, align environment, bracket-matrix
\usepackage{amsmath}
\usepackage{amssymb}
%% allow page breaks within display mathematics anywhere
%% level 4 is maximally permissive
%% this is exactly the opposite of AMSmath package philosophy
%% there are per-display, and per-equation options to control this
%% split, aligned, gathered, and alignedat are not affected
\allowdisplaybreaks[4]
%% allow more columns to a matrix
%% can make this even bigger by overriding with  latex.preamble.late  processing option
\setcounter{MaxMatrixCols}{30}
%%
%%
%% Division Titles, and Page Headers/Footers
%% titlesec package, loading "titleps" package cooperatively
%% See code comments about the necessity and purpose of "explicit" option
\usepackage[explicit, pagestyles]{titlesec}
\newtitlemark{\chaptertitlename}
%% Set global/default page style for document due
%% to potential re-definitions after documentclass
\pagestyle{headings}
%%
%% Create globally-available macros to be provided for style writers
%% These are redefined for each occurence of each division
\newcommand{\divisionnameptx}{\relax}%
\newcommand{\titleptx}{\relax}%
\newcommand{\subtitleptx}{\relax}%
\newcommand{\shortitleptx}{\relax}%
\newcommand{\authorsptx}{\relax}%
\newcommand{\epigraphptx}{\relax}%
%% Create environments for possible occurences of each division
%% Environment for a PTX "preface" at the level of a LaTeX "chapter"
\NewDocumentEnvironment{preface}{mmmmmm}
{%
\renewcommand{\divisionnameptx}{Preface}%
\renewcommand{\titleptx}{#1}%
\renewcommand{\subtitleptx}{#2}%
\renewcommand{\shortitleptx}{#3}%
\renewcommand{\authorsptx}{#4}%
\renewcommand{\epigraphptx}{#5}%
\chapter*{#1}%
\addcontentsline{toc}{chapter}{#3}
\label{#6}%
}{}%
%% Environment for a PTX "part" at the level of a LaTeX "part"
\NewDocumentEnvironment{partptx}{mmmmmm}
{%
\renewcommand{\divisionnameptx}{Part}%
\renewcommand{\titleptx}{#1}%
\renewcommand{\subtitleptx}{#2}%
\renewcommand{\shortitleptx}{#3}%
\renewcommand{\authorsptx}{#4}%
\renewcommand{\epigraphptx}{#5}%
\part[#3]{#1}%
\label{#6}%
}{}%
%% Environment for a PTX "chapter" at the level of a LaTeX "chapter"
\NewDocumentEnvironment{chapterptx}{mmmmmm}
{%
\renewcommand{\divisionnameptx}{Chapter}%
\renewcommand{\titleptx}{#1}%
\renewcommand{\subtitleptx}{#2}%
\renewcommand{\shortitleptx}{#3}%
\renewcommand{\authorsptx}{#4}%
\renewcommand{\epigraphptx}{#5}%
\chapter[#3]{#1}%
\label{#6}%
}{}%
%% Environment for a PTX "section" at the level of a LaTeX "section"
\NewDocumentEnvironment{sectionptx}{mmmmmm}
{%
\renewcommand{\divisionnameptx}{Section}%
\renewcommand{\titleptx}{#1}%
\renewcommand{\subtitleptx}{#2}%
\renewcommand{\shortitleptx}{#3}%
\renewcommand{\authorsptx}{#4}%
\renewcommand{\epigraphptx}{#5}%
\section[#3]{#1}%
\label{#6}%
}{}%
%% Environment for a PTX "reading-questions" at the level of a LaTeX "subsection"
\NewDocumentEnvironment{reading-questions-subsection}{mmmmmm}
{%
\renewcommand{\divisionnameptx}{Reading Questions}%
\renewcommand{\titleptx}{#1}%
\renewcommand{\subtitleptx}{#2}%
\renewcommand{\shortitleptx}{#3}%
\renewcommand{\authorsptx}{#4}%
\renewcommand{\epigraphptx}{#5}%
\subsection[#3]{#1}%
\label{#6}%
}{}%
%% Environment for a PTX "reading-questions" at the level of a LaTeX "subsection"
\NewDocumentEnvironment{reading-questions-subsection-numberless}{mmmmmm}
{%
\renewcommand{\divisionnameptx}{Reading Questions}%
\renewcommand{\titleptx}{#1}%
\renewcommand{\subtitleptx}{#2}%
\renewcommand{\shortitleptx}{#3}%
\renewcommand{\authorsptx}{#4}%
\renewcommand{\epigraphptx}{#5}%
\subsection*{#1}%
\addcontentsline{toc}{subsection}{#3}
\label{#6}%
}{}%
%% Environment for a PTX "subsection" at the level of a LaTeX "subsection"
\NewDocumentEnvironment{subsectionptx}{mmmmmm}
{%
\renewcommand{\divisionnameptx}{Subsection}%
\renewcommand{\titleptx}{#1}%
\renewcommand{\subtitleptx}{#2}%
\renewcommand{\shortitleptx}{#3}%
\renewcommand{\authorsptx}{#4}%
\renewcommand{\epigraphptx}{#5}%
\subsection[#3]{#1}%
\label{#6}%
}{}%
%%
%% Styles for six traditional LaTeX divisions
\titleformat{\chapter}[display]
{\normalfont\huge\bfseries}{\divisionnameptx\space\thechapter}{20pt}{\Huge#1}
[{\Large\authorsptx}]
\titleformat{name=\chapter,numberless}[display]
{\normalfont\huge\bfseries}{}{0pt}{#1}
[{\Large\authorsptx}]
\titlespacing*{\chapter}{0pt}{50pt}{40pt}
\titleformat{\section}[hang]
{\normalfont\Large\bfseries}{\thesection}{1ex}{#1}
[{\large\authorsptx}]
\titleformat{name=\section,numberless}[block]
{\normalfont\Large\bfseries}{}{0pt}{#1}
[{\large\authorsptx}]
\titlespacing*{\section}{0pt}{3.5ex plus 1ex minus .2ex}{2.3ex plus .2ex}
\titleformat{\subsection}[hang]
{\normalfont\large\bfseries}{\thesubsection}{1ex}{#1}
[{\normalsize\authorsptx}]
\titleformat{name=\subsection,numberless}[block]
{\normalfont\large\bfseries}{}{0pt}{#1}
[{\normalsize\authorsptx}]
\titlespacing*{\subsection}{0pt}{3.25ex plus 1ex minus .2ex}{1.5ex plus .2ex}
\titleformat{\subsubsection}[hang]
{\normalfont\normalsize\bfseries}{\thesubsubsection}{1em}{#1}
[{\small\authorsptx}]
\titleformat{name=\subsubsection,numberless}[block]
{\normalfont\normalsize\bfseries}{}{0pt}{#1}
[{\normalsize\authorsptx}]
\titlespacing*{\subsubsection}{0pt}{3.25ex plus 1ex minus .2ex}{1.5ex plus .2ex}
\titleformat{\paragraph}[hang]
{\normalfont\normalsize\bfseries}{\theparagraph}{1em}{#1}
[{\small\authorsptx}]
\titleformat{name=\paragraph,numberless}[block]
{\normalfont\normalsize\bfseries}{}{0pt}{#1}
[{\normalsize\authorsptx}]
\titlespacing*{\paragraph}{0pt}{3.25ex plus 1ex minus .2ex}{1.5em}
%%
%% Semantic Macros
%% To preserve meaning in a LaTeX file
%% Only defined here if required in this document
%% Used for warnings, typically bold and italic
\newcommand{\alert}[1]{\textbf{\textit{#1}}}
%% Used for inline definitions of terms
\newcommand{\terminology}[1]{\textbf{#1}}
%% Used for fillin answer blank
%% Argument is length in em
%% Length may compress for output to fit in one line
\newcommand{\fillin}[1]{\leavevmode\leaders\vrule height -1.2pt depth 1.5pt \hskip #1em minus #1em \null}
%% Division Numbering: Chapters, Sections, Subsections, etc
%% Division numbers may be turned off at some level ("depth")
%% A section *always* has depth 1, contrary to us counting from the document root
%% The latex default is 3.  If a larger number is present here, then
%% removing this command may make some cross-references ambiguous
%% The precursor variable $numbering-maxlevel is checked for consistency in the common XSL file
\setcounter{secnumdepth}{3}
%% begin: General AMS environment setup
%% Environments built with amsthm package
\usepackage{amsthm}
%% Numbering for Theorems, Conjectures, Examples, Figures, etc
%% Controlled by  numbering.theorems.level  processing parameter
%% Numbering: all theorem-like numbered consecutively
%% i.e. Corollary 4.3 follows Theorem 4.2
%% Always need some theorem environment to set base numbering scheme
%% even if document has no theorems (but has other environments)
%% Create a never-used style first, always
%% simply to provide a global counter to use, namely "cthm"
\newtheorem{cthm}{BadTheoremStringName}[section]
%% AMS "proof" environment is not used, but we leave previously
%% implemented \qedhere in place, should the LaTeX be recycled
\renewcommand{\qedhere}{\relax}
%% end: General AMS environment setup
%%
%% tcolorbox, with styles, for DEFINITION-LIKE
%%
%% definition: fairly simple numbered block/structure
\tcbset{ definitionstyle/.style={bwminimalstyle, runintitlestyle, blockspacingstyle, after title={\space}, after upper={\space\space\hspace*{\stretch{1}}\(\lozenge\)}, } }
\newtcolorbox[use counter*=cthm]{definition}[2]{title={{Definition~\thecthm\notblank{#1}{\space\space#1}{}}}, phantomlabel={#2}, breakable, parbox=false, definitionstyle, }
%%
%% tcolorbox, with styles, for COMPUTATION-LIKE
%%
%% computation: fairly simple numbered block/structure
\tcbset{ computationstyle/.style={bwminimalstyle, runintitlestyle, blockspacingstyle, after title={\space}, } }
\newtcolorbox[use counter*=cthm]{computation}[2]{title={{Computation~\thecthm\notblank{#1}{\space\space#1}{}}}, phantomlabel={#2}, breakable, parbox=false, computationstyle, }
%%
%% tcolorbox, with styles, for EXAMPLE-LIKE
%%
%% example: fairly simple numbered block/structure
\tcbset{ examplestyle/.style={bwminimalstyle, runintitlestyle, blockspacingstyle, after title={\space}, after upper={\space\space\hspace*{\stretch{1}}\(\square\)}, } }
\newtcolorbox[use counter*=cthm]{example}[2]{title={{Example~\thecthm\notblank{#1}{\space\space#1}{}}}, phantomlabel={#2}, breakable, parbox=false, examplestyle, }
%% Divisional exercises (and worksheet) as LaTeX environments
%% Third argument is option for extra workspace in worksheets
%% Hanging indent occupies a 5ex width slot prior to left margin
%% Experimentally this seems just barely sufficient for a bold "888."
%% Division exercises, not in exercise group
\tcbset{ divisionexercisestyle/.style={bwminimalstyle, runintitlestyle, exercisespacingstyle, left=5ex, breakable, parbox=false } }
\newtcolorbox{divisionexercise}[4]{divisionexercisestyle, before title={\hspace{-5ex}\makebox[5ex][l]{#1.}}, title={\notblank{#2}{#2\space}{}}, phantom={\hypertarget{#4}{}}, after={\notblank{#3}{\newline\rule{\workspacestrutwidth}{#3\textheight}\newline}{}}}
%% Localize LaTeX supplied names (possibly none)
\renewcommand*{\partname}{Part}
\renewcommand*{\chaptername}{Chapter}
%% Equation Numbering
%% Controlled by  numbering.equations.level  processing parameter
\numberwithin{equation}{section}
%% For improved tables
\usepackage{array}
%% Some extra height on each row is desirable, especially with horizontal rules
%% Increment determined experimentally
\setlength{\extrarowheight}{0.2ex}
%% Define variable thickness horizontal rules, full and partial
%% Thicknesses are 0.03, 0.05, 0.08 in the  booktabs  package
\newcommand{\hrulethin}  {\noalign{\hrule height 0.04em}}
\newcommand{\hrulemedium}{\noalign{\hrule height 0.07em}}
\newcommand{\hrulethick} {\noalign{\hrule height 0.11em}}
%% We preserve a copy of the \setlength package before other
%% packages (extpfeil) get a chance to load packages that redefine it
\let\oldsetlength\setlength
\newlength{\Oldarrayrulewidth}
\newcommand{\crulethin}[1]%
{\noalign{\global\oldsetlength{\Oldarrayrulewidth}{\arrayrulewidth}}%
\noalign{\global\oldsetlength{\arrayrulewidth}{0.04em}}\cline{#1}%
\noalign{\global\oldsetlength{\arrayrulewidth}{\Oldarrayrulewidth}}}%
\newcommand{\crulemedium}[1]%
{\noalign{\global\oldsetlength{\Oldarrayrulewidth}{\arrayrulewidth}}%
\noalign{\global\oldsetlength{\arrayrulewidth}{0.07em}}\cline{#1}%
\noalign{\global\oldsetlength{\arrayrulewidth}{\Oldarrayrulewidth}}}
\newcommand{\crulethick}[1]%
{\noalign{\global\oldsetlength{\Oldarrayrulewidth}{\arrayrulewidth}}%
\noalign{\global\oldsetlength{\arrayrulewidth}{0.11em}}\cline{#1}%
\noalign{\global\oldsetlength{\arrayrulewidth}{\Oldarrayrulewidth}}}
%% Single letter column specifiers defined via array package
\newcolumntype{A}{!{\vrule width 0.04em}}
\newcolumntype{B}{!{\vrule width 0.07em}}
\newcolumntype{C}{!{\vrule width 0.11em}}
%% Figures, Tables, Listings, Named Lists, Floats
%% The [H]ere option of the float package fixes floats in-place,
%% in deference to web usage, where floats are totally irrelevant
%% You can remove some of this setup, to restore standard LaTeX behavior
%% HOWEVER, numbering of figures/tables AND theorems/examples/remarks, etc
%% may de-synchronize with the numbering in the HTML version
%% You can remove the "placement={H}" option to allow flotation and
%% preserve numbering, BUT the numbering may then appear "out-of-order"
%% Floating environments: http://tex.stackexchange.com/questions/95631/
\usepackage{float}
\usepackage{newfloat}
\usepackage{caption}%% Adjust stock figure environment so that it no longer floats
\SetupFloatingEnvironment{figure}{fileext=lof,placement={H},within=section,name=Figure}
\captionsetup[figure]{labelfont=bf}
%% http://tex.stackexchange.com/questions/16195
\makeatletter
\let\c@figure\c@cthm
\makeatother
%% Adjust stock table environment so that it no longer floats
\SetupFloatingEnvironment{table}{fileext=lot,placement={H},within=section,name=Table}
\captionsetup[table]{labelfont=bf}
%% http://tex.stackexchange.com/questions/16195
\makeatletter
\let\c@table\c@cthm
\makeatother
%% QR Code Support
%% Videos and other interactives
\usepackage{qrcode}
\newlength{\qrsize}
\newlength{\previewwidth}
%% tcolorbox styles for interactive previews
%% changing size= and/or colback can aid in debugging
\tcbset{ previewstyle/.style={bwminimalstyle, halign=center} }
\tcbset{ qrstyle/.style={bwminimalstyle, hbox} }
\tcbset{ captionstyle/.style={bwminimalstyle, left=1em, width=\linewidth} }
%% Generic red play button (from SVG)
%% tikz package should be loaded by now
\definecolor{playred}{RGB}{230,33,23}
\newcommand{\genericpreview}{
        \begin{tikzpicture}[y=0.80pt, x=0.80pt, yscale=-1.000000, xscale=1.000000, inner sep=0pt, outer sep=0pt]
        \path[fill=playred] (94.9800,28.8400) .. controls (94.9800,28.8400) and
        (94.0400,22.2400) .. (91.1700,19.3400) .. controls (87.5300,15.5300) and
        (83.4500,15.5100) .. (81.5800,15.2900) .. controls (68.1800,14.3200) and
        (48.0600,14.4400) .. (48.0600,14.4400) .. controls (48.0600,14.4400) and
        (27.9400,14.3200) .. (14.5400,15.2900) .. controls (12.6700,15.5100) and
        (8.5900,15.5300) .. (4.9500,19.3400) .. controls (2.0800,22.2400) and
        (1.1400,28.8400) .. (1.1400,28.8400) .. controls (1.1400,28.8400) and
        (0.1800,36.5800) .. (0.0000,44.3300) -- (0.0000,51.5900) .. controls
        (0.1800,59.3400) and (1.1400,67.0800) .. (1.1400,67.0800) .. controls
        (1.1400,67.0800) and (2.0700,73.6800) .. (4.9500,76.5800) .. controls
        (8.5900,80.3900) and (13.3800,80.2700) .. (15.5100,80.6700) .. controls
        (23.0400,81.3900) and (47.2100,81.5600) .. (48.0500,81.5700) .. controls
        (48.0600,81.5700) and (68.1900,81.6000) .. (81.5900,80.6300) .. controls
        (83.4600,80.4100) and (87.5400,80.3900) .. (91.1800,76.5800) .. controls
        (94.0500,73.6800) and (94.9900,67.0800) .. (94.9900,67.0800) .. controls
        (94.9900,67.0800) and (95.9500,59.3300) .. (96.0100,51.5900) --
        (96.0100,44.3300) .. controls (95.9400,36.5800) and (94.9800,28.8400) ..
        (94.9800,28.8400) -- cycle(38.2800,61.4100) -- (38.2800,34.4100) --
        (64.0200,47.9100) -- (38.2800,61.4100) -- cycle;
        \end{tikzpicture}
        }
%% Fancy Verbatim for consoles, preformatted, code display, literate programming
\usepackage{fancyvrb}
%% code display (cd), by analogy with math display (md)
%% savebox, lrbox, etc to achieve centering
\newsavebox{\codedisplaybox}
\newenvironment{codedisplay}
{\VerbatimEnvironment\begin{center}\begin{lrbox}{\codedisplaybox}\begin{BVerbatim}}
{\end{BVerbatim}\end{lrbox}\usebox{\codedisplaybox}\end{center}}
%% More flexible list management, esp. for references
%% But also for specifying labels (i.e. custom order) on nested lists
\usepackage{enumitem}
%% hyperref driver does not need to be specified, it will be detected
%% Footnote marks in tcolorbox have broken linking under
%% hyperref, so it is necessary to turn off all linking
%% It *must* be given as a package option, not with \hypersetup
\usepackage[hyperfootnotes=false]{hyperref}
%% Hyperlinking active in electronic PDFs, all links solid and blue
\hypersetup{colorlinks=true,linkcolor=blue,citecolor=blue,filecolor=blue,urlcolor=blue}
\hypersetup{pdftitle={Statistics Lecture Notes}}
%% If you manually remove hyperref, leave in this next command
\providecommand\phantomsection{}
%% If tikz has been loaded, replace ampersand with \amp macro
%% tcolorbox styles for sidebyside layout
\tcbset{ sbsstyle/.style={raster equal height=rows,raster force size=false} }
\tcbset{ sbsheadingstyle/.style={bwminimalstyle, halign=center, fontupper=\bfseries} }
\tcbset{ sbspanelstyle/.style={bwminimalstyle} }
\tcbset{ sbscaptionstyle/.style={bwminimalstyle, halign=center} }
%% Enviroments for side-by-side and components
%% Necessary to use \NewTColorBox for boxes of the panels
%% "newfloat" environment to squash page-breaks within a single sidebyside
%% \leavevmode necessary when a side-by-side comes first, right after a heading
%% "xparse" environment for entire sidebyside
\NewDocumentEnvironment{sidebyside}{mmmm}
  {\begin{tcbraster}
    [sbsstyle,raster columns=#1,
    raster left skip=#2\linewidth,raster right skip=#3\linewidth,raster column skip=#4\linewidth]}
  {\end{tcbraster}}
%% "tcolorbox" environments for three components of a panel
\NewTColorBox{sbsheading}{m}{sbsheadingstyle,width=#1\linewidth}
\NewTColorBox{sbspanel}{mO{top}}{sbspanelstyle,width=#1\linewidth,valign=#2}
\NewTColorBox{sbscaption}{m}{sbscaptionstyle,width=#1\linewidth}
%% extpfeil package for certain extensible arrows,
%% as also provided by MathJax extension of the same name
%% NB: this package loads mtools, which loads calc, which redefines
%%     \setlength, so it can be removed if it seems to be in the 
%%     way and your math does not use:
%%     
%%     \xtwoheadrightarrow, \xtwoheadleftarrow, \xmapsto, \xlongequal, \xtofrom
%%     
%%     we have had to be extra careful with variable thickness
%%     lines in tables, and so also load this package late
\usepackage{extpfeil}
%% Custom Preamble Entries, late (use latex.preamble.late)
%% Begin: Author-provided packages
%% (From  docinfo/latex-preamble/package  elements)
%% End: Author-provided packages
%% Begin: Author-provided macros
%% (From  docinfo/macros  element)
%% Plus three from MBX for XML characters

\newcommand{\lt}{<}
\newcommand{\gt}{>}
\newcommand{\amp}{&}
%% End: Author-provided macros
\begin{document}
\frontmatter
%% begin: half-title
\thispagestyle{empty}
{\centering
\vspace*{0.28\textheight}
{\Huge Statistics Lecture Notes}\\}
\clearpage
%% end:   half-title
%% begin: adcard
\thispagestyle{empty}
\null%
\clearpage
%% end:   adcard
%% begin: title page
%% Inspired by Peter Wilson's "titleDB" in "titlepages" CTAN package
\thispagestyle{empty}
{\centering
\vspace*{0.14\textheight}
%% Target for xref to top-level element is ToC
\addtocontents{toc}{\protect\hypertarget{index}{}}
{\Huge Statistics Lecture Notes}\\[3\baselineskip]
{\Large Jason D. Underdown}\\[0.5\baselineskip]
{\Large SLCC}\\[3\baselineskip]
{\Large May 29, 2019 22:02:12 (-06:00)}\\}
\clearpage
%% end:   title page
%% begin: copyright-page
\thispagestyle{empty}
\hypertarget{colophon-1}{}\vspace*{\stretch{2}}
\noindent{\bfseries Website}: \href{http:\slash{}\slash{}www.math.utah.edu\slash{}\textasciitilde{}jasonu\slash{}stats\slash{}}{Statistics Lecture Notes}\par\medskip
\noindent\textcopyright{}2019\quad{}Jason D. Underdown\par\medskip
\vspace*{\stretch{1}}
\null\clearpage
%% end:   copyright-page
%
%
\typeout{************************************************}
\typeout{Preface  Preface}
\typeout{************************************************}
%
\begin{preface}{Preface}{}{Preface}{}{}{preface-1}
\hypertarget{p-1}{}%
As the course progresses, I will add the example problems that I present in class at the chalkboard to these lecture notes.%
\par
\hypertarget{p-2}{}%
Studying these lecture notes alone is usually not sufficient to achieve an A grade in the course. You should prepare for each lecture by reading the section we will cover each day before coming to class. Do not feel you need to read each section extremely carefully. If you get stuck while reading, spend a few minutes trying to figure out your misunderstanding, but if you can't figure it out, make a mental note of your question (or better yet write it down) and then continue on in the section. Often when reading mathematics, you will find that your question gets answered a few sentences or paragraphs later.%
\end{preface}
%% begin: table of contents
%% Adjust Table of Contents
\setcounter{tocdepth}{0}
\renewcommand*\contentsname{Contents}
\tableofcontents
%% end:   table of contents
\mainmatter
%
%
\typeout{************************************************}
\typeout{Part I Getting the Information You Need}
\typeout{************************************************}
%
\begin{partptx}{Getting the Information You Need}{}{Getting the Information You Need}{}{}{part-getting-information}
 %
%
\typeout{************************************************}
\typeout{Chapter 1 Data Collection}
\typeout{************************************************}
%
\begin{chapterptx}{Data Collection}{}{Data Collection}{}{}{ch-data-collection}
%
%
\typeout{************************************************}
\typeout{Section 1.1 Introduction to the Practice of Statistics}
\typeout{************************************************}
%
\begin{sectionptx}{Introduction to the Practice of Statistics}{}{Introduction to the Practice of Statistics}{}{}{sec-practice-of-stats}
\begin{definition}{statistics.}{definition-1}%
how to collect, organize, summarize and analyze information so that conclusions can be drawn with a measure of confidence.\end{definition}
\begin{definition}{descriptive stats.}{definition-2}%
entails organizing and summarizing collected data.\end{definition}
\begin{definition}{inferential stats.}{definition-3}%
methods for taking sample results and extending it to the population with some measure of reliability.\end{definition}
\begin{definition}{population.}{definition-4}%
entire group of interest in a study\end{definition}
\begin{definition}{individual.}{definition-5}%
a single person or object being studied\end{definition}
\begin{definition}{sample.}{definition-6}%
subset of the population being studied\end{definition}
\begin{definition}{statistic.}{definition-7}%
numerical summary of a sample\end{definition}
\begin{definition}{parameter.}{definition-8}%
numerical summary of a population\end{definition}
\begin{definition}{variables.}{definition-9}%
characteristics of individuals in the population\end{definition}
\begin{definition}{qualitative variable.}{definition-10}%
responses based on attributes or characteristics\end{definition}
\begin{definition}{quantitative variable.}{definition-11}%
responses are numerical measures, with a natural ordering i.e.\@ it makes sense to compare via \textless{} or \textgreater{}\end{definition}
\begin{definition}{discrete variable.}{definition-12}%
finite or countable number of possible values e.g.\@ number of children\end{definition}
\begin{definition}{continuous variable.}{definition-13}%
infinite number of possible values e.g.\@ height\end{definition}
\begin{figure}
\centering
\setlength{\qrsize}{9em}
\setlength{\previewwidth}{\linewidth}
\addtolength{\previewwidth}{-\qrsize}
\begin{tcbraster}[raster columns=2, raster column skip=1pt, raster halign=center, raster force size=false, raster left skip=0pt, raster right skip=0pt]%
\begin{tcolorbox}[previewstyle, width=\previewwidth]%
\includegraphics[width=0.80\linewidth,height=\qrsize,keepaspectratio]{images/video-1.jpg}%
\end{tcolorbox}%
\begin{tcolorbox}[qrstyle]%
{\hypersetup{urlcolor=black}\qrcode[height=\qrsize]{https://www.youtube.com/watch?v=sxQaBpKfDRk}}%
\end{tcolorbox}%
\begin{tcolorbox}[captionstyle]%
\small YouTube: \mono{https://www.youtube.com/watch?v=sxQaBpKfDRk}\end{tcolorbox}%
\end{tcbraster}%
\caption{Crash Course Statistics: What Is Statistics\label{figure-1}}
\end{figure}
%
%
\typeout{************************************************}
\typeout{Reading Questions 1.1 Statistic or Parameter?}
\typeout{************************************************}
%
\begin{reading-questions-subsection-numberless}{Statistic or Parameter?}{}{Statistic or Parameter?}{}{}{reading-questions-1}
\hypertarget{p-3}{}%
For each description which follows, determine whether the numerical quantity is a \emph{statistic} or \emph{parameter}.%
\begin{divisionexercise}{1}{}{}{exercise-1}%
In a study of all 2223 passengers aboard the Titanic, it is found that 706 survived when it sank. \textbf{Answer}.\hypertarget{answer-1}{}\quad%
parameter\end{divisionexercise}%
\begin{divisionexercise}{2}{}{}{exercise-2}%
A recent survey of a sample of MBAs reported that the average salary for an MBA is more than \textdollar{}82,000. (Source: The Wall Street Journal) \textbf{Answer}.\hypertarget{answer-2}{}\quad%
statistic\end{divisionexercise}%
\begin{divisionexercise}{3}{}{}{exercise-3}%
Starting salaries for all of the 667 MBA graduates from the University of Chicago Graduate School of Business increased 8.5\% from the previous year. \textbf{Answer}.\hypertarget{answer-3}{}\quad%
parameter\end{divisionexercise}%
\begin{divisionexercise}{4}{}{}{exercise-4}%
In a recent poll of Salt Lake Community College students, 83\% of students owned a vehicle. \textbf{Answer}.\hypertarget{answer-4}{}\quad%
statistic\end{divisionexercise}%
\begin{divisionexercise}{5}{}{}{exercise-5}%
In a recent survey of 457 attendees of Jumanji, 402 would recommend the movie to a friend. \textbf{Answer}.\hypertarget{answer-5}{}\quad%
parameter, unless you consider the population to be all people who have seen Jumanji\end{divisionexercise}%
\end{reading-questions-subsection-numberless}
%
%
\typeout{************************************************}
\typeout{Reading Questions 1.1 Discrete or Continuous Variable?}
\typeout{************************************************}
%
\begin{reading-questions-subsection-numberless}{Discrete or Continuous Variable?}{}{Discrete or Continuous Variable?}{}{}{reading-questions-2}
\begin{divisionexercise}{1}{}{}{exercise-6}%
In a survey of 1059 adults, it is found that 39\% of them have guns in their homes \textbf{Answer}.\hypertarget{answer-6}{}\quad%
discrete\end{divisionexercise}%
\begin{divisionexercise}{2}{}{}{exercise-7}%
The number of heads obtained after flipping a coin five times \textbf{Answer}.\hypertarget{answer-7}{}\quad%
discrete\end{divisionexercise}%
\begin{divisionexercise}{3}{}{}{exercise-8}%
The distance Tiger Woods can drive a golf ball \textbf{Answer}.\hypertarget{answer-8}{}\quad%
continuous\end{divisionexercise}%
\begin{divisionexercise}{4}{}{}{exercise-9}%
Points scored in a basketball game \textbf{Answer}.\hypertarget{answer-9}{}\quad%
discrete\end{divisionexercise}%
\begin{divisionexercise}{5}{}{}{exercise-10}%
Volume of water lost each day from a leaky faucet \textbf{Answer}.\hypertarget{answer-10}{}\quad%
continuous\end{divisionexercise}%
\begin{divisionexercise}{6}{}{}{exercise-11}%
Length of a song \textbf{Answer}.\hypertarget{answer-11}{}\quad%
continuous, but could be discrete if measured in seconds\end{divisionexercise}%
\begin{divisionexercise}{7}{}{}{exercise-12}%
Number of words in a song \textbf{Answer}.\hypertarget{answer-12}{}\quad%
discrete\end{divisionexercise}%
\end{reading-questions-subsection-numberless}
%
%
\typeout{************************************************}
\typeout{Reading Questions 1.1 Qualitative or Quantitative Variable?}
\typeout{************************************************}
%
\begin{reading-questions-subsection-numberless}{Qualitative or Quantitative Variable?}{}{Qualitative or Quantitative Variable?}{}{}{reading-questions-3}
\begin{divisionexercise}{1}{}{}{exercise-13}%
Gender \textbf{Answer}.\hypertarget{answer-13}{}\quad%
qualitative\end{divisionexercise}%
\begin{divisionexercise}{2}{}{}{exercise-14}%
Temperature \textbf{Answer}.\hypertarget{answer-14}{}\quad%
quantitative\end{divisionexercise}%
\begin{divisionexercise}{3}{}{}{exercise-15}%
Nation of origin \textbf{Answer}.\hypertarget{answer-15}{}\quad%
qualitative\end{divisionexercise}%
\begin{divisionexercise}{4}{}{}{exercise-16}%
Number of siblings \textbf{Answer}.\hypertarget{answer-16}{}\quad%
quantitative\end{divisionexercise}%
\begin{divisionexercise}{5}{}{}{exercise-17}%
Number of days \textbf{Answer}.\hypertarget{answer-17}{}\quad%
quantitative\end{divisionexercise}%
\begin{divisionexercise}{6}{}{}{exercise-18}%
Grams of carbohydrates in a donut \textbf{Answer}.\hypertarget{answer-18}{}\quad%
quantitative\end{divisionexercise}%
\begin{divisionexercise}{7}{}{}{exercise-19}%
Phone number \textbf{Answer}.\hypertarget{answer-19}{}\quad%
qualitative\end{divisionexercise}%
\begin{divisionexercise}{8}{}{}{exercise-20}%
Value of a house \textbf{Answer}.\hypertarget{answer-20}{}\quad%
quantitative\end{divisionexercise}%
\begin{divisionexercise}{9}{}{}{exercise-21}%
Zip code \textbf{Answer}.\hypertarget{answer-21}{}\quad%
qualitative\end{divisionexercise}%
\end{reading-questions-subsection-numberless}
%
%
\typeout{************************************************}
\typeout{Reading Questions 1.1 2011 BMW Models}
\typeout{************************************************}
%
\begin{reading-questions-subsection-numberless}{2011 BMW Models}{}{2011 BMW Models}{}{}{reading-questions-4}
\hypertarget{p-4}{}%
The following information relates to the 2011 model year product line of BMW automobiles. Identify the individuals in the study, the variables and if the data corresponding to the variables are qualitative or quantitative, and continuous or discrete.%
\begin{table}
\centering
\begin{tabular}{llll}\hrulethick
Model&Body Style&Weight&Number of Seats\tabularnewline\hrulemedium
3 Series&Coupe&\multicolumn{1}{r}{3362}&\multicolumn{1}{c}{4}\tabularnewline[0pt]
5 Series&Sedan&\multicolumn{1}{r}{4056}&\multicolumn{1}{c}{5}\tabularnewline[0pt]
6 Series&Convertible&\multicolumn{1}{r}{4277}&\multicolumn{1}{c}{4}\tabularnewline[0pt]
7 Series&Sedan&\multicolumn{1}{r}{4564}&\multicolumn{1}{c}{5}\tabularnewline[0pt]
X3&Sport Utility&\multicolumn{1}{r}{4012}&\multicolumn{1}{c}{5}\tabularnewline[0pt]
Z4&Coupe&\multicolumn{1}{r}{3505}&\multicolumn{1}{c}{2}\tabularnewline\hrulethick
\end{tabular}
\caption{2011 BMW Automobiles\label{table-1}}
\end{table}
\begin{divisionexercise}{1}{}{}{exercise-22}%
Individuals: \textbf{Answer}.\hypertarget{answer-22}{}\quad%
3 Series, 5 Series, 6 Series, 7 Series, X3, Z4\end{divisionexercise}%
\begin{divisionexercise}{2}{}{}{exercise-23}%
Variable(s): \textbf{Answer}.\hypertarget{answer-23}{}\quad%
\begin{tabular}{lll}\hrulethick
Variable&Qual.\slash{}Quant.&Disc.\slash{}Cont.\tabularnewline\hrulemedium
body style&qualitative&N\slash{}A\tabularnewline[0pt]
weight&quantitative&continuous\tabularnewline[0pt]
number of seats&quantitative&discrete\tabularnewline\hrulethick
\end{tabular}
\end{divisionexercise}%
\end{reading-questions-subsection-numberless}
%
%
\typeout{************************************************}
\typeout{Reading Questions 1.1 Smoker's IQ}
\typeout{************************************************}
%
\begin{reading-questions-subsection-numberless}{Smoker's IQ}{}{Smoker's IQ}{}{}{reading-questions-5}
\hypertarget{p-5}{}%
A study was conducted in which 20,211 18-{}-{}year old recruits were given an exam to measure IQ. In addition, the recruits were asked to disclose their smoking status. An individual was considered a smoker if he smoked at least 1 cigarette per day. The goal of the study was to determine whether adolescents aged 18 to 21 who smoke have a lower IQ than nonsmokers. It was found that the average IQ of the smokers was 94, while the average IQ of nonsmokers was 101. The researchers concluded that lower IQ individuals are more likely to chose to smoke, not that smoking makes people less intelligent.%
\begin{divisionexercise}{1}{}{}{exercise-24}%
What is the population being studied? \textbf{Answer}.\hypertarget{answer-24}{}\quad%
adolescents aged 18 to 21\end{divisionexercise}%
\begin{divisionexercise}{2}{}{}{exercise-25}%
What is the sample? \textbf{Answer}.\hypertarget{answer-25}{}\quad%
20,211 18 year old recruits\end{divisionexercise}%
\begin{divisionexercise}{3}{}{}{exercise-26}%
What are the descriptive statistics? \textbf{Answer}.\hypertarget{answer-26}{}\quad%
IQ, smoking status\end{divisionexercise}%
\end{reading-questions-subsection-numberless}
\begin{figure}
\centering
\setlength{\qrsize}{9em}
\setlength{\previewwidth}{\linewidth}
\addtolength{\previewwidth}{-\qrsize}
\begin{tcbraster}[raster columns=2, raster column skip=1pt, raster halign=center, raster force size=false, raster left skip=0pt, raster right skip=0pt]%
\begin{tcolorbox}[previewstyle, width=\previewwidth]%
\includegraphics[width=0.80\linewidth,height=\qrsize,keepaspectratio]{images/video-2.jpg}%
\end{tcolorbox}%
\begin{tcolorbox}[qrstyle]%
{\hypersetup{urlcolor=black}\qrcode[height=\qrsize]{https://www.youtube.com/watch?v=tN9Xl1AcSv8}}%
\end{tcolorbox}%
\begin{tcolorbox}[captionstyle]%
\small YouTube: \mono{https://www.youtube.com/watch?v=tN9Xl1AcSv8}\end{tcolorbox}%
\end{tcbraster}%
\caption{Crash Course Statistics: Mathematical Thinking\label{figure-2}}
\end{figure}
\end{sectionptx}
%
%
\typeout{************************************************}
\typeout{Section 1.2 Observational Studies Vs. Designed Experiments}
\typeout{************************************************}
%
\begin{sectionptx}{Observational Studies Vs. Designed Experiments}{}{Observational Studies Vs. Designed Experiments}{}{}{sec-observational-vs-designed}
\begin{definition}{response variable.}{definition-14}%
The response variable is the focus of a question in a study or experiment.\end{definition}
\begin{definition}{explanatory variable.}{definition-15}%
An explanatory variable is one that explains changes in the response variable.\end{definition}
\begin{definition}{observational study.}{definition-16}%
measuring response variable without influencing the values of either the response variable or explanatory variable\end{definition}
\begin{definition}{cross-sectional studies.}{definition-17}%
observational studies that collect data at one particular time\end{definition}
\begin{definition}{case-control studies.}{definition-18}%
observational studies that require individuals to look back in time to record measurements (retrospective)\end{definition}
\begin{definition}{cohort studies.}{definition-19}%
observational studies that track individuals repeatedly over time (prospective)\end{definition}
\begin{definition}{designed experiment.}{definition-20}%
measuring response variable after intentionally changing the value of the explanatory variable\end{definition}
\begin{definition}{confounding.}{definition-21}%
in a study occurs when the effects of two or more explanatory variables are not separated. Therefore, any relation that may exist between an explanatory variable and the response variable may be due to some other variable or variables not accounted for in the study\end{definition}
\begin{definition}{lurking variable.}{definition-22}%
a variable that was not considered in the study, but affects the response variable\end{definition}
\hypertarget{p-6}{}%
\alert{IMPORTANT: Because of the possibility of lurking variables, observational studies \emph{DO NOT} allow a researcher to claim causation, only \emph{association}.}%
\begin{definition}{confounding variable.}{definition-23}%
a variable in the study where you can’t tell the difference between its effect and another explanatory variable’s effect\end{definition}
\begin{definition}{census.}{definition-24}%
list of information from all individuals in a population\end{definition}
\begin{figure}
\centering
\setlength{\qrsize}{9em}
\setlength{\previewwidth}{\linewidth}
\addtolength{\previewwidth}{-\qrsize}
\begin{tcbraster}[raster columns=2, raster column skip=1pt, raster halign=center, raster force size=false, raster left skip=0pt, raster right skip=0pt]%
\begin{tcolorbox}[previewstyle, width=\previewwidth]%
\includegraphics[width=0.80\linewidth,height=\qrsize,keepaspectratio]{images/video-3.jpg}%
\end{tcolorbox}%
\begin{tcolorbox}[qrstyle]%
{\hypersetup{urlcolor=black}\qrcode[height=\qrsize]{https://www.youtube.com/watch?v=kkBDa-ICvyY}}%
\end{tcolorbox}%
\begin{tcolorbox}[captionstyle]%
\small YouTube: \mono{https://www.youtube.com/watch?v=kkBDa-ICvyY}\end{tcolorbox}%
\end{tcbraster}%
\caption{Crash Course Statistics: Controlled Experiments\label{figure-3}}
\end{figure}
%
%
\typeout{************************************************}
\typeout{Reading Questions 1.2 Observational Study or Designed Experiment?}
\typeout{************************************************}
%
\begin{reading-questions-subsection-numberless}{Observational Study or Designed Experiment?}{}{Observational Study or Designed Experiment?}{}{}{reading-questions-6}
\begin{divisionexercise}{1}{}{}{exercise-27}%
Rats with cancer are divided into two groups. One group receives 5 milligrams of a medication that is thought to fight cancer, and the other receives 10 milligrams. After 2 years, the spread of cancer is measured. \textbf{Answer}.\hypertarget{answer-27}{}\quad%
designed experiment\end{divisionexercise}%
\begin{divisionexercise}{2}{}{}{exercise-28}%
Conservation agents netted 250 large-mouth bass in a lake and determined how many were carrying parasites. \textbf{Answer}.\hypertarget{answer-28}{}\quad%
observational study\end{divisionexercise}%
\begin{divisionexercise}{3}{}{}{exercise-29}%
Seventh grade students are randomly divided into two groups. One group is taught math using traditional techniques; the other is taught math using a reform method. After 1 year, each group is given an achievement test to compare proficiency. \textbf{Answer}.\hypertarget{answer-29}{}\quad%
designed experiment\end{divisionexercise}%
\begin{divisionexercise}{4}{}{}{exercise-30}%
A survey was conducted asking 400 people, "Do you prefer Coke or Pepsi?" \textbf{Answer}.\hypertarget{answer-30}{}\quad%
observational study\end{divisionexercise}%
\end{reading-questions-subsection-numberless}
%
%
\typeout{************************************************}
\typeout{Reading Questions 1.2 Daily Coffee Consumption}
\typeout{************************************************}
%
\begin{reading-questions-subsection-numberless}{Daily Coffee Consumption}{}{Daily Coffee Consumption}{}{}{reading-questions-7}
\hypertarget{p-7}{}%
Is there an association between daily coffee consumption and the occurrence of skin cancer? Reasearchers asked 93,676 women to disclose their coffee drinking habits and also determined which of the women had nonmelanoma skin cancer. The researchers concluded that consumption of 6 or more cups of caffeinated coffee per day was associated with a reduction in nonmelanoma skin cancer.%
\begin{divisionexercise}{1}{}{}{exercise-31}%
Was this an oberservational study or experiment? If observational, what type of observational study was it? \textbf{Answer}.\hypertarget{answer-31}{}\quad%
observational, cross-section\end{divisionexercise}%
\begin{divisionexercise}{2}{}{}{exercise-32}%
What is the response variable? What is the explanatory variable? \textbf{Answer}.\hypertarget{answer-32}{}\quad%
response: skin cancer, explanatory: coffee consumption\end{divisionexercise}%
\begin{divisionexercise}{3}{}{}{exercise-33}%
Can we conclude drinking six or more cups of coffee reduces the chance of nonmelanoma skin cancer? \textbf{Answer}.\hypertarget{answer-33}{}\quad%
No, because observational studies never allow us to determine causation, just association.\end{divisionexercise}%
\end{reading-questions-subsection-numberless}
%
%
\typeout{************************************************}
\typeout{Reading Questions 1.2 Get Married, Gain Weight}
\typeout{************************************************}
%
\begin{reading-questions-subsection-numberless}{Get Married, Gain Weight}{}{Get Married, Gain Weight}{}{}{reading-questions-8}
\hypertarget{p-8}{}%
Are young couples who marry or cohabitate more likely to gain weight than those who stay single? Researchers followed 8000 men and women for 7 years. At the start of the study, none of the participants were married or living with a romantic partner. The researchers found that women who married or cohabitated during the study period gained 9 pounds more then single women, and married or cohabitating men gained, on average, 6 pounds more than single men.%
\begin{divisionexercise}{1}{}{}{exercise-34}%
Why is this an observational study? What type of observational study is it? \textbf{Answer}.\hypertarget{answer-34}{}\quad%
No manipulation of explanatory variable; cohort study\end{divisionexercise}%
\begin{divisionexercise}{2}{}{}{exercise-35}%
What is the response variable in the study? What is the explanatory variable? \textbf{Answer}.\hypertarget{answer-35}{}\quad%
response: weight, explanatory: marital status\end{divisionexercise}%
\begin{divisionexercise}{3}{}{}{exercise-36}%
Identify some of the potential lurking variables in the study. \textbf{Answer}.\hypertarget{answer-36}{}\quad%
diet, exercise, stress, pregnancy\end{divisionexercise}%
\begin{divisionexercise}{4}{}{}{exercise-37}%
Can we conclude that getting married or cohabitating causes one to gain weight? \textbf{Answer}.\hypertarget{answer-37}{}\quad%
No, because it is an observational study, not an experiment. We can only conclude that they are associated.\end{divisionexercise}%
\end{reading-questions-subsection-numberless}
\end{sectionptx}
%
%
\typeout{************************************************}
\typeout{Section 1.3 Simple Random Sampling}
\typeout{************************************************}
%
\begin{sectionptx}{Simple Random Sampling}{}{Simple Random Sampling}{}{}{sec-simple-random-sampling}
\begin{definition}{random sampling.}{definition-25}%
using chance (an objective device) to select individuals from a population to be included in the sample\end{definition}
\begin{definition}{simple random sample.}{definition-26}%
every sample of a certain size is equally likely (and every subject has an equal chance of being selected)\end{definition}
\begin{definition}{\(N\).}{definition-27}%
population size\end{definition}
\begin{definition}{\(n\).}{definition-28}%
sample size\end{definition}
\begin{definition}{frame.}{definition-29}%
list of all individuals in a population being studied\end{definition}
\begin{definition}{sampling without replacement.}{definition-30}%
once selected, an individual cannot be chosen again\end{definition}
\begin{definition}{sampling with replacement.}{definition-31}%
once selected, individuals are placed back in the population and can be chosen again\end{definition}
%
%
\typeout{************************************************}
\typeout{Subsection 1.3.1 Generating Random Numbers on the TI-84 Plus Calculator}
\typeout{************************************************}
%
\begin{subsectionptx}{Generating Random Numbers on the TI-84 Plus Calculator}{}{Generating Random Numbers on the TI-84 Plus Calculator}{}{}{sub-generating-random-numbers}
\hypertarget{p-9}{}%
There are several ways to generate random numbers on your calculator. One simple way is to use the \mono{randInt} function.%
\begin{computation}{Generate Random Integers.}{computation-1}%
\hypertarget{p-10}{}%
To generate a random integer between 1 and 10:\leavevmode%
\begin{enumerate}
\item\hypertarget{li-1}{}Press the \mono{math} button and use the cursor to highlight the \mono{PROB} tab.%
\item\hypertarget{li-2}{}Select \mono{5:randInt(}%
\item\hypertarget{li-3}{}Fill out the StatWizard so it looks like below%
\begin{codedisplay}

		lower:1
		upper:10
		n:
		Paste
	      
\end{codedisplay}
Highlight \mono{Paste} and press \mono{enter}.%
\item\hypertarget{li-4}{}Press \mono{enter} to execute the function you previously pasted into the command window.%
\end{enumerate}
%
\begin{figure}\centering\includegraphics[width=0.4\linewidth]{images/randInt-single.png}
\end{figure}\end{computation}
\begin{computation}{List of Random Integers.}{computation-2}%
\hypertarget{p-11}{}%
To generate a list of 4 random integers between 1 and 10:\leavevmode%
\begin{enumerate}
\item\hypertarget{li-5}{}Press the \mono{math} button and use the cursor to highlight the \mono{PROB} tab.%
\item\hypertarget{li-6}{}Select \mono{5:randInt(}%
\item\hypertarget{li-7}{}Fill out the StatWizard so it looks like below%
\begin{codedisplay}

		lower:1
		upper:10
		n:4
		Paste
	      
\end{codedisplay}
Highlight \mono{Paste} and press \mono{enter}.%
\item\hypertarget{li-8}{}Press \mono{enter} to execute the function you previously pasted into the command window.%
\end{enumerate}
%
\begin{figure}\centering\includegraphics[width=0.4\linewidth]{images/randInt-list.png}
\end{figure}\end{computation}
\begin{computation}{List of Random Integers with No Repeats.}{computation-3}%
\hypertarget{p-12}{}%
To generate a list of 10 random integers between 1 and 20 with \emph{no repeats}:\leavevmode%
\begin{enumerate}
\item\hypertarget{li-9}{}Press the \mono{math} button and use the cursor to highlight the \mono{PROB} tab.%
\item\hypertarget{li-10}{}Select \mono{8:randIntNoRep(}%
\item\hypertarget{li-11}{}Fill out the StatWizard so it looks like below%
\begin{codedisplay}

		lower:1
		upper:20
		n:10
		Paste
	      
\end{codedisplay}
Highlight \mono{Paste} and press \mono{enter}.%
\item\hypertarget{li-12}{}Press \mono{enter} to execute the function you previously pasted into the command window.%
\end{enumerate}
%
\begin{figure}\centering\includegraphics[width=0.4\linewidth]{images/randIntNoRep.png}
\end{figure}\end{computation}
\end{subsectionptx}
%
%
\typeout{************************************************}
\typeout{Reading Questions 1.3.2 Four Tickets for Friends}
\typeout{************************************************}
%
\begin{reading-questions-subsection}{Four Tickets for Friends}{}{Four Tickets for Friends}{}{}{reading-questions-9}
\hypertarget{p-13}{}%
Sophia has 10 friends that she would like to invite to a concert, but only 4 tickets to give away.%
\par
\hypertarget{p-14}{}%
Mike, Jamie, Adam, Yvette, Ashley, Monica, Cherie, Julie, Willard, Bruce%
\begin{divisionexercise}{1}{}{}{exercise-38}%
Which of the following would produce a simple random sample of the 4 friends that she will bring with her: \leavevmode%
\begin{enumerate}[label=(\alph*)]
\item\hypertarget{li-13}{}List each persons name on a piece of paper, place them in a hat and draw 4. \textbf{Answer}.\quad%
yes%
\item\hypertarget{li-14}{}List the names in alphabetical order and take the first 4 names. \textbf{Answer}.\quad%
no%
\item\hypertarget{li-15}{}Ask one of her friends who she should bring. \textbf{Answer}.\quad%
no%
\item\hypertarget{li-16}{}Number the friends from 1 to 10 and use a random number generator to produce 4 numbers between 1 and 10 which correspond to the 10 friends. \textbf{Answer}.\quad%
yes%
\end{enumerate}
\end{divisionexercise}%
\begin{divisionexercise}{2}{}{}{exercise-39}%
Use the TI graphing calculator to obtain a random sample of the 4 friends that she will bring with her and give their names. \textbf{Solution}.\hypertarget{solution-1}{}\quad%
\begin{figure}\centering\includegraphics[width=0.4\linewidth]{images/four-friends.png}
\end{figure}\hypertarget{p-15}{}%
Thus she would invite Mike, Yvette, Ashley and Julie.%
\par
\hypertarget{p-16}{}%
We can also have the calculator sort the list in ascending order:%
\begin{figure}\centering\includegraphics[width=0.4\linewidth]{images/four-friends-sorted.png}
\end{figure}\hypertarget{p-17}{}%
In this case, she would invite Mike, Jamie, Monica and Willard.%
\end{divisionexercise}%
\end{reading-questions-subsection}
%
%
\typeout{************************************************}
\typeout{Reading Questions 1.3.3 Employee Survey}
\typeout{************************************************}
%
\begin{reading-questions-subsection}{Employee Survey}{}{Employee Survey}{}{}{reading-questions-10}
\hypertarget{p-18}{}%
The owner of a private food store is concerned with employee morale. She decides to survey the employees to learn about work environment and job satisfaction. \begin{tabular}{llll}\hrulethick
1 Archer&9 Foushi&16 Kemp&23 Oliver\tabularnewline\hrulethin
2 Bolcerek&10 Gow&17 Lathus&24 Orsini\tabularnewline\hrulethin
3 Bryant&11 Grove&18 Lindsey&25 Salazar\tabularnewline\hrulethin
4 Carlisle&12 Hall&19 Massie&26 Ullrich\tabularnewline\hrulethin
5 Cole&13 Hills&20 McGuffin&27 Vaneck\tabularnewline\hrulethin
6 Dimas&14 Houston&21 Musa&28 Weber\tabularnewline\hrulethin
7 Ellison&15 Kats&22 Nickas&29 Zavodny\tabularnewline\hrulethin
8 Everhart&&&\tabularnewline\hrulethick
\end{tabular}
%
\begin{divisionexercise}{1}{}{}{exercise-40}%
Obtain a simple random sample of size 6 from the above table using the TI graphing calculator. \textbf{Solution}.\hypertarget{solution-2}{}\quad%
\begin{figure}\centering\includegraphics[width=0.4\linewidth]{images/employee-survey.png}
\end{figure} Thus the randomly selected employees are: Bolcerek, Carlisle, Musa, Vaneck, Weber, and Zavodny.\end{divisionexercise}%
\end{reading-questions-subsection}
\end{sectionptx}
%
%
\typeout{************************************************}
\typeout{Section 1.4 Other Effective Sampling Methods}
\typeout{************************************************}
%
\begin{sectionptx}{Other Effective Sampling Methods}{}{Other Effective Sampling Methods}{}{}{sec-other-sampling-methods}
\begin{definition}{stratified sampling.}{definition-32}%
a simple random sample is drawn from each nonoverlapping subgroup (or \terminology{stratum}) into which the population has been separated.\hypertarget{p-19}{}%
Works best when individuals in each stratum are similar in some way and we want to make sure each stratum type is represented. The size of each simple random sample is proportional to the number of individuals in the stratum from which it is selected. (e.g.\@ separating individuals into age groups)%
\end{definition}
\begin{definition}{systematic sampling.}{definition-33}%
\hypertarget{p-20}{}%
a starting point, \(p\), is randomly selected and then every \(k\)th subject is included in the sample, where%
\begin{gather*}
k = \lfloor N/n \rfloor
\end{gather*}
The brackets, \(\lfloor x \rfloor\), mean round \(x\) down to the nearest integer, or take the \terminology{floor} of \(x\). Subjects in the sample are numbered:%
\begin{gather*}
p, p+k, p+2k, p+3k, \ldots, p+(n-1)k
\end{gather*}
%
\par
\hypertarget{p-21}{}%
Works best when individuals are arranged sequentially but you’re not sure how many total individuals there are because you don’t have a frame of the population (e.g.\@ files in a drawer, people coming out of a building exit polling, etc.\@).%
\end{definition}
\begin{definition}{cluster sampling.}{definition-34}%
\hypertarget{p-22}{}%
the population is divided into sections or \terminology{clusters}; all individuals from randomly selected clusters are included in the sample.%
\par
\hypertarget{p-23}{}%
Works best when individuals are grouped geographically so it makes logistical sense to sample this way, but location does not affect your variable of interest; clusters are heterogeneous like the population.%
\end{definition}
\begin{definition}{convenience sampling.}{definition-35}%
\hypertarget{p-24}{}%
sampling subjects who are easy to get (bad sampling method).%
\par
\hypertarget{p-25}{}%
Works best when: NEVER!! (not random)%
\end{definition}
\begin{definition}{self-selected sampling.}{definition-36}%
\hypertarget{p-26}{}%
a convenience sample where individuals are self-selected (another bad sampling method).%
\par
\hypertarget{p-27}{}%
Works best when: NEVER!! (not random, favors those with stronger opinions)%
\end{definition}
\begin{definition}{multi-stage sampling.}{definition-37}%
\hypertarget{p-28}{}%
a sample that is selected in multiple stages, each of which might use a different method of sampling.%
\par
\hypertarget{p-29}{}%
Works best when: The population is organized hierarchically but no list exists.%
\end{definition}
%
%
\typeout{************************************************}
\typeout{Reading Questions 1.4 Identify the Sampling Method Part I}
\typeout{************************************************}
%
\begin{reading-questions-subsection-numberless}{Identify the Sampling Method Part I}{}{Identify the Sampling Method Part I}{}{}{reading-questions-11}
\begin{divisionexercise}{1}{}{}{exercise-41}%
The fifth person exiting the student center and every 12th person afterward is interviewed. \textbf{Answer}.\hypertarget{answer-42}{}\quad%
systematic sampling\end{divisionexercise}%
\begin{divisionexercise}{2}{}{}{exercise-42}%
Two out of the six parking lots on campus are randomly selected and every driver in the two lots receives a questionnaire. \textbf{Answer}.\hypertarget{answer-43}{}\quad%
cluster sampling\end{divisionexercise}%
\begin{divisionexercise}{3}{}{}{exercise-43}%
Each registered student is assigned a number, 20 numbers are randomly generated and the corresponding students are interviewed. \textbf{Answer}.\hypertarget{answer-44}{}\quad%
simple random sampling\end{divisionexercise}%
\begin{divisionexercise}{4}{}{}{exercise-44}%
Six buildings on campus are randomly selected and every 4th person who exits each building is included in the sample. \textbf{Answer}.\hypertarget{answer-45}{}\quad%
multi-stage sampling; (cluster sampling followed by systematic sampling)\end{divisionexercise}%
\begin{divisionexercise}{5}{}{}{exercise-45}%
Randomly select 10 males and 10 females in the activity center to be interviewed. \textbf{Answer}.\hypertarget{answer-46}{}\quad%
stratified sampling\end{divisionexercise}%
\begin{divisionexercise}{6}{}{}{exercise-46}%
Every person in my class that sits next to me in Math 1040 is interviewed. \textbf{Answer}.\hypertarget{answer-47}{}\quad%
convenience sampling\end{divisionexercise}%
\begin{divisionexercise}{7}{}{}{exercise-47}%
A pile of questionnaires is left on a table in the library with a sign that says, “Please complete.” \textbf{Answer}.\hypertarget{answer-48}{}\quad%
self-selected sampling\end{divisionexercise}%
\end{reading-questions-subsection-numberless}
%
%
\typeout{************************************************}
\typeout{Reading Questions 1.4 Identify the Sampling Method Part II}
\typeout{************************************************}
%
\begin{reading-questions-subsection-numberless}{Identify the Sampling Method Part II}{}{Identify the Sampling Method Part II}{}{}{reading-questions-12}
\hypertarget{p-30}{}%
For each scenario determine which type of sampling is used.%
\begin{divisionexercise}{1}{}{}{exercise-48}%
To estimate the percentage of defects in a recent manufacturing batch, a quality control manager selects every 8th chip that comes of the assemply line starting with the 3rd until she obtains 140 chips. \textbf{Answer}.\hypertarget{answer-49}{}\quad%
systematic sampling\end{divisionexercise}%
\begin{divisionexercise}{2}{}{}{exercise-49}%
To determine customer opinion of its boarding policy, Southwest Airlines randomly selects 60 flights during a certain week and surveys all passengers on the flights. \textbf{Answer}.\hypertarget{answer-50}{}\quad%
cluster sampling\end{divisionexercise}%
\begin{divisionexercise}{3}{}{}{exercise-50}%
A member of Congress wishes to determine her constituency’s opinions regarding estate taxes. She first divides her constituency into three income classes: low-income households, middle-income households, and upper-income households. She then takes a simple random sample of households from each income class. \textbf{Answer}.\hypertarget{answer-51}{}\quad%
stratified sampling\end{divisionexercise}%
\begin{divisionexercise}{4}{}{}{exercise-51}%
In an effort to identify whether an advertising campaign has been effective, a marketing firm conducts a nation wide poll by randomly selecting individuals from a list of known users of the product. \textbf{Answer}.\hypertarget{answer-52}{}\quad%
simple random sampling\end{divisionexercise}%
\begin{divisionexercise}{5}{}{}{exercise-52}%
A radio station asks its listeners to call in their opinion regarding the use of U.S. forces in peace keeping missions. \textbf{Answer}.\hypertarget{answer-53}{}\quad%
self-selected sampling\end{divisionexercise}%
\begin{divisionexercise}{6}{}{}{exercise-53}%
A college official divides the student population into five classes: freshman, sophomore, junior, senior and graduate student. The official takes a simple random sample from each class and asks the members’ opinions regarding student services. \textbf{Answer}.\hypertarget{answer-54}{}\quad%
stratified sampling\end{divisionexercise}%
\end{reading-questions-subsection-numberless}
%
%
\typeout{************************************************}
\typeout{Reading Questions 1.4 Systematic Sampling Example}
\typeout{************************************************}
%
\begin{reading-questions-subsection-numberless}{Systematic Sampling Example}{}{Systematic Sampling Example}{}{}{reading-questions-13}
\hypertarget{p-31}{}%
The human resources department at a company wants to conduct a survey regarding employee morale. The department has a list of all 400 employees at the company and wants to do a systematic sample of size 20. To do this they randomly select the 8th employee and then select every 20th person. Which workers will be surveyed?%
\begin{divisionexercise}{1}{}{}{exercise-54}%
1st = \fillin{4.545454545454546}\textbf{Answer}.\hypertarget{answer-55}{}\quad%
\(8 = p\)\end{divisionexercise}%
\begin{divisionexercise}{2}{}{}{exercise-55}%
2nd = \fillin{4.545454545454546}\textbf{Answer}.\hypertarget{answer-56}{}\quad%
\(28 = p + k = 8 + 20\)\end{divisionexercise}%
\begin{divisionexercise}{3}{}{}{exercise-56}%
3rd = \fillin{4.545454545454546}\textbf{Answer}.\hypertarget{answer-57}{}\quad%
\(48 = p + 2k = 8 + 40\)\end{divisionexercise}%
\begin{divisionexercise}{4}{}{}{exercise-57}%
20th = \fillin{4.545454545454546}\textbf{Answer}.\hypertarget{answer-58}{}\quad%
\(388 = p + (n-1)k = 8 + (19)20\)\end{divisionexercise}%
\end{reading-questions-subsection-numberless}
\end{sectionptx}
%
%
\typeout{************************************************}
\typeout{Section 1.5 Bias in Sampling}
\typeout{************************************************}
%
\begin{sectionptx}{Bias in Sampling}{}{Bias in Sampling}{}{}{sec-bias-in-sampling}
\begin{definition}{bias.}{definition-38}%
when sample results are not representative of the population\end{definition}
\begin{definition}{sample bias.}{definition-39}%
bias because the technique used to select the sample favors some individuals over others (not random)\end{definition}
\begin{definition}{undercoverage.}{definition-40}%
type of sampling bias where part of the population has a lower chance or no chance of being in the sample (e.g.\@ if the frame is incomplete)\end{definition}
\begin{definition}{nonresponse bias.}{definition-41}%
bias because individuals in the sample do not respond\end{definition}
\begin{definition}{response bias.}{definition-42}%
bias because responses do not reflect the true feelings or beliefs of the respondent.~\hypertarget{p-32}{}%
\terminology{Response bias} can occur in a number of different ways:\leavevmode%
\begin{itemize}[label=\textbullet]
\item{}From the interviewer\textemdash{}the nature or wording of the interviewer can affect the truthfulness of responses.%
\item{}Misrepresented answers\textemdash{}when the results of a survey misrepresent facts or are flat out lies.%
\item{}Wording of questions\textemdash{}the way a question is worded can lead to response bias in a survey.%
\item{}Ordering of questions or choices\textemdash{}many surveys will rearrange the order of the questions so that responses are not influenced by prior answers.%
\end{itemize}
%
\end{definition}
\begin{definition}{open question.}{definition-43}%
allows a respondent to choose his or her own response e.g.\@ "What is the most important problem facing youth today?"\end{definition}
\begin{definition}{closed question.}{definition-44}%
requires the respondent to choose from a list of predetermined responses.\end{definition}
\begin{definition}{nonsampling errors.}{definition-45}%
errors that result from undercoverage, non-response bias, response bias or data-entry error.\end{definition}
\begin{definition}{sampling errors.}{definition-46}%
errors that result from using a sample to estimate information about a population.\end{definition}
\hypertarget{p-33}{}%
I will leave it for you to read section 1.5 for examples. This section is only three pages but is a must read section!%
\begin{figure}
\centering
\setlength{\qrsize}{9em}
\setlength{\previewwidth}{\linewidth}
\addtolength{\previewwidth}{-\qrsize}
\begin{tcbraster}[raster columns=2, raster column skip=1pt, raster halign=center, raster force size=false, raster left skip=0pt, raster right skip=0pt]%
\begin{tcolorbox}[previewstyle, width=\previewwidth]%
\includegraphics[width=0.80\linewidth,height=\qrsize,keepaspectratio]{images/video-4.jpg}%
\end{tcolorbox}%
\begin{tcolorbox}[qrstyle]%
{\hypersetup{urlcolor=black}\qrcode[height=\qrsize]{https://www.youtube.com/watch?v=Rf-fIpB4D50}}%
\end{tcolorbox}%
\begin{tcolorbox}[captionstyle]%
\small YouTube: \mono{https://www.youtube.com/watch?v=Rf-fIpB4D50}\end{tcolorbox}%
\end{tcbraster}%
\caption{Crash Course Statistics: Sampling Methods and Bias in Surveys\label{figure-4}}
\end{figure}
\end{sectionptx}
\end{chapterptx}
\end{partptx}
%
%
\typeout{************************************************}
\typeout{Part II Descriptive Statistics}
\typeout{************************************************}
%
\begin{partptx}{Descriptive Statistics}{}{Descriptive Statistics}{}{}{part-descriptive-stats}
 %
%
\typeout{************************************************}
\typeout{Chapter 2 Summarizing Data in Tables and Graphs}
\typeout{************************************************}
%
\begin{chapterptx}{Summarizing Data in Tables and Graphs}{}{Summarizing Data in Tables and Graphs}{}{}{ch-tables-and-graphs}
%
%
\typeout{************************************************}
\typeout{Section 2.1 Organizing Qualitative Data}
\typeout{************************************************}
%
\begin{sectionptx}{Organizing Qualitative Data}{}{Organizing Qualitative Data}{}{}{sec-organizing-qualitative-data}
\begin{definition}{frequency distribution.}{definition-47}%
\hypertarget{p-34}{}%
lists each category of data and the number of occurrences for each category.%
\end{definition}
\begin{definition}{relative frequency distribution.}{definition-48}%
\hypertarget{p-35}{}%
lists each category of data and the proportion (or percent) of obervations in that category.%
\par
\hypertarget{p-36}{}%
relative frequency = count\slash{}total%
\end{definition}
\begin{definition}{pie chart.}{definition-49}%
\hypertarget{p-37}{}%
is a circle divided into sectors. Each sector represents a category of data and its area is proportional to the frequency.%
\end{definition}
\begin{figure}
\centering
\setlength{\qrsize}{9em}
\setlength{\previewwidth}{\linewidth}
\addtolength{\previewwidth}{-\qrsize}
\begin{tcbraster}[raster columns=2, raster column skip=1pt, raster halign=center, raster force size=false, raster left skip=0pt, raster right skip=0pt]%
\begin{tcolorbox}[previewstyle, width=\previewwidth]%
\includegraphics[width=0.80\linewidth,height=\qrsize,keepaspectratio]{images/video-5.jpg}%
\end{tcolorbox}%
\begin{tcolorbox}[qrstyle]%
{\hypersetup{urlcolor=black}\qrcode[height=\qrsize]{https://www.youtube.com/watch?v=hEWY6kkBdpo}}%
\end{tcolorbox}%
\begin{tcolorbox}[captionstyle]%
\small YouTube: \mono{https://www.youtube.com/watch?v=hEWY6kkBdpo}\end{tcolorbox}%
\end{tcbraster}%
\caption{Crash Course Statistics: Data Visualization Part 1\label{figure-5}}
\end{figure}
%
%
\typeout{************************************************}
\typeout{Reading Questions 2.1 Favorite Day to Eat Out}
\typeout{************************************************}
%
\begin{reading-questions-subsection-numberless}{Favorite Day to Eat Out}{}{Favorite Day to Eat Out}{}{}{reading-questions-14}
\hypertarget{p-38}{}%
A survey was conducted by Wakefield Research in which 40 particpants were asked to give their favorite night to go out to dinner. The following data are the results.%
\begin{table}
\centering
\begin{tabular}{lllll}\hrulethick
Thursday&Saturday&Friday&Friday&Sunday\tabularnewline\hrulethin
Wednesday&Saturday&Friday&Tuesday&Friday\tabularnewline\hrulethin
Saturday&Monday&Friday&Friday&Sunday\tabularnewline\hrulethin
Friday&Tuesday&Wednesday&Saturday&Friday\tabularnewline\hrulethin
Wednesday&Monday&Wednesday&Wednesday&Friday\tabularnewline\hrulethin
Friday&Wednesday&Thursday&Tuesday&Friday\tabularnewline\hrulethin
Tuesday&Saturday&Friday&Tuesday&Friday\tabularnewline\hrulethin
Saturday&Saturday&Saturday&Sunday&Friday\tabularnewline\hrulethin
\end{tabular}
\caption{Data from Wakefield Research\label{table-2}}
\end{table}
\begin{divisionexercise}{1}{}{}{exercise-58}%
\hypertarget{p-39}{}%
Construct a frequency distribution%
\begin{tabular}{lll}\hrulethick
Day&Tally&Frequency\tabularnewline\hrulethick
Sunday& ~~~~~ ~~~~~ ~~~~~ ~~~~~ ~~~~~&\tabularnewline\hrulethin
Monday&&\tabularnewline\hrulethin
Tuesday&&\tabularnewline\hrulethin
Wednesday&&\tabularnewline\hrulethin
Thursday&&\tabularnewline\hrulethin
Friday&&\tabularnewline\hrulethin
Saturday&&\tabularnewline\hrulethick
Sum:&&\tabularnewline\hrulethick
\end{tabular}
\textbf{Solution}.\hypertarget{solution-3}{}\quad%
\begin{tabular}{ll}\hrulethick
Day&\multicolumn{1}{c}{Frequency}\tabularnewline\hrulethick
Sunday&\multicolumn{1}{c}{3}\tabularnewline\hrulethin
Monday&\multicolumn{1}{c}{2}\tabularnewline\hrulethin
Tuesday&\multicolumn{1}{c}{5}\tabularnewline\hrulethin
Wednesday&\multicolumn{1}{c}{6}\tabularnewline\hrulethin
Thursday&\multicolumn{1}{c}{2}\tabularnewline\hrulethin
Friday&\multicolumn{1}{c}{14}\tabularnewline\hrulethin
Saturday&\multicolumn{1}{c}{8}\tabularnewline\hrulethick
Sum:&\multicolumn{1}{c}{40}\tabularnewline\hrulethick
\end{tabular}
\end{divisionexercise}%
\begin{divisionexercise}{2}{}{}{exercise-59}%
\hypertarget{p-40}{}%
Construct a relative frequency distribution%
\begin{tabular}{ll}\hrulethick
Day&\multicolumn{1}{c}{Relative Frequency}\tabularnewline\hrulethick
Sunday&\multicolumn{1}{c}{}\tabularnewline\hrulethin
Monday&\multicolumn{1}{c}{}\tabularnewline\hrulethin
Tuesday&\multicolumn{1}{c}{}\tabularnewline\hrulethin
Wednesday&\multicolumn{1}{c}{}\tabularnewline\hrulethin
Thursday&\multicolumn{1}{c}{}\tabularnewline\hrulethin
Friday&\multicolumn{1}{c}{}\tabularnewline\hrulethin
Saturday&\multicolumn{1}{c}{}\tabularnewline\hrulethick
Sum:&\multicolumn{1}{c}{}\tabularnewline\hrulethick
\end{tabular}
\textbf{Solution}.\hypertarget{solution-4}{}\quad%
\begin{tabular}{ll}\hrulethick
Day&\multicolumn{1}{r}{Relative Frequency}\tabularnewline\hrulethick
Sunday&\multicolumn{1}{r}{0.075}\tabularnewline\hrulethin
Monday&\multicolumn{1}{r}{0.050}\tabularnewline\hrulethin
Tuesday&\multicolumn{1}{r}{0.125}\tabularnewline\hrulethin
Wednesday&\multicolumn{1}{r}{0.150}\tabularnewline\hrulethin
Thursday&\multicolumn{1}{r}{0.050}\tabularnewline\hrulethin
Friday&\multicolumn{1}{r}{0.350}\tabularnewline\hrulethin
Saturday&\multicolumn{1}{r}{0.200}\tabularnewline\hrulethick
Sum:&\multicolumn{1}{r}{1.000}\tabularnewline\hrulethick
\end{tabular}
\end{divisionexercise}%
\begin{divisionexercise}{3}{}{}{exercise-60}%
\hypertarget{p-41}{}%
Construct a frequency bar chart%
\textbf{Solution}.\hypertarget{solution-5}{}\quad%
\begin{figure}\centering\includegraphics[width=1\linewidth]{images/eat-out-bar-chart.png}
\end{figure}\end{divisionexercise}%
\begin{divisionexercise}{4}{}{}{exercise-61}%
\hypertarget{p-42}{}%
Construct a relative frequency bar chart%
\textbf{Solution}.\hypertarget{solution-6}{}\quad%
\begin{figure}\centering\includegraphics[width=1\linewidth]{images/eat-out-rel-freq-bar-chart.png}
\end{figure}\end{divisionexercise}%
\begin{divisionexercise}{5}{}{}{exercise-62}%
\hypertarget{p-43}{}%
Construct a pie chart%
\textbf{Solution}.\hypertarget{solution-7}{}\quad%
\begin{figure}\centering\includegraphics[width=1\linewidth]{images/eat-out-pie-chart.png}
\end{figure}\end{divisionexercise}%
\begin{divisionexercise}{6}{}{}{exercise-63}%
\hypertarget{p-44}{}%
If you own a restaurant, which days would you purchase an advertisement on the local radio? Are there days that you should avoid?%
\textbf{Answer}.\hypertarget{answer-59}{}\quad%
\hypertarget{p-45}{}%
It would be good to advertise on Fridays and Saturdays, but not so worthwhile to advertise on Mondays and Thursdays.%
\end{divisionexercise}%
\end{reading-questions-subsection-numberless}
\end{sectionptx}
%
%
\typeout{************************************************}
\typeout{Section 2.2 Organizing Quantitative Data}
\typeout{************************************************}
%
\begin{sectionptx}{Organizing Quantitative Data}{}{Organizing Quantitative Data}{}{}{sec-organizing-quantitative-data}
\begin{definition}{}{definition-50}%
\hypertarget{p-46}{}%
A \terminology{bar graph} is for qualitative data while a \terminology{histogram} is used for quantitative data. In a histogram, data is grouped into \terminology{classes} or \terminology{bins}.%
\par
\hypertarget{p-47}{}%
Notes on histograms:\leavevmode%
\begin{enumerate}
\item\hypertarget{li-21}{}Bars have equal widths%
\item\hypertarget{li-22}{}Bars are touching%
\end{enumerate}
%
\par
\hypertarget{p-48}{}%
Terminology related to histograms:\leavevmode%
\begin{itemize}[label=\textbullet]
\item{}\terminology{lower class limit}: smallest value in a class%
\item{}\terminology{upper class limit}: largest value in a class%
\item{}\terminology{class width (bin width)}: difference between consecutive lower class limits (or upper class limits)%
\item{}\terminology{open ended}: a table where the first class has no lower limit and\slash{}or the last class has no upper limit.%
\end{itemize}
%
\end{definition}
%
%
\typeout{************************************************}
\typeout{Subsection 2.2.1 The Shape of a Distribution}
\typeout{************************************************}
%
\begin{subsectionptx}{The Shape of a Distribution}{}{The Shape of a Distribution}{}{}{sub-distribution-shapes}
\hypertarget{p-49}{}%
One reason for going through the trouble of displaying data in a histogram is to determine the basic shape of your data. The four most common shapes are:%
\begin{figure}
\centering
\includegraphics[width=0.8\linewidth]{images/four-distribution-shapes.png}
\caption{Four most common distribution shapes.\label{figure-6}}
\end{figure}
\end{subsectionptx}
%
%
\typeout{************************************************}
\typeout{Subsection 2.2.2 Create Histograms and Frequency Distributions from Data Using the TI-84 Plus}
\typeout{************************************************}
%
\begin{subsectionptx}{Create Histograms and Frequency Distributions from Data Using the TI-84 Plus}{}{Create Histograms and Frequency Distributions from Data Using the TI-84 Plus}{}{}{sub-histogram-from-data}
\hypertarget{p-50}{}%
The TI-84 Plus calculator can generate histograms and frequency distributions from raw data. Normally, one would first sort their data, separate it into classes to create a frequency distribution and then plot the frequency distribution with vertical bars to generate a histogram. However, when we use the calculator, we will first generate the histogram and then read off the frequencies from the height of each bar to create the frequency distribution.%
\par
\hypertarget{p-51}{}%
The data below was generated by asking each student from a statistics class to give their height in inches.%
\begin{table}
\centering
\begin{tabular}{llllll}\hrulethick
62&79&63&68&62&69\tabularnewline\hrulethin
69&64&62&63&68&71\tabularnewline\hrulethin
72&68&63&70&63&67\tabularnewline\hrulethin
60&66&71&69&70&71\tabularnewline\hrulethick
\end{tabular}
\caption{Student heights (inches)\label{tab-student-heights}}
\end{table}
\hypertarget{p-52}{}%
To create a histogram of the height data above do the following:%
\leavevmode%
\begin{enumerate}
\item\hypertarget{li-27}{}Enter the data in the above table into list 1. This can be done by pressing \mono{stat}, \mono{1} and then using the cursor to select list 1.%
\item\hypertarget{li-28}{}\hypertarget{p-53}{}%
Press \mono{2nd}, \mono{y=} to enter the STAT PLOTS menu. (Make sure Plot2 and Plot3 are OFF by pressing \mono{4} (PlotsOff) if necessary.) Choose Plot1 by pressing \mono{1}.%
\begin{figure}\centering\includegraphics[width=0.4\linewidth]{images/stat-plots-menu.png}
\end{figure}\item\hypertarget{li-29}{}\hypertarget{p-54}{}%
Make your screen match the image below:%
\begin{figure}\centering\includegraphics[width=0.4\linewidth]{images/stat-plot-histogram-from-list.png}
\end{figure}\item\hypertarget{li-30}{}\hypertarget{p-55}{}%
Press \mono{zoom}, \mono{9}. This will cause the calculator to bin the data according to its internal algorithm and display the histogram.%
\begin{figure}\centering\includegraphics[width=0.4\linewidth]{images/heights-histogram-1.png}
\end{figure}\hypertarget{p-56}{}%
You can press \mono{trace} and move the cursor to check the class\slash{}bin limits the calculator chose.%
\item\hypertarget{li-31}{}\hypertarget{p-57}{}%
Press \mono{window} and modify Xmin, Xmax, Xscl, Ymin and Ymax to match the screen below. Setting Xscl to 2 changes the class\slash{}bin width to 2 (inches).%
\begin{figure}\centering\includegraphics[width=0.4\linewidth]{images/heights-histogram-window.png}
\end{figure}\item\hypertarget{li-32}{}\hypertarget{p-58}{}%
Press \mono{graph} to display the new histogram.%
\begin{figure}\centering\includegraphics[width=0.4\linewidth]{images/heights-histogram-2.png}
\end{figure}\hypertarget{p-59}{}%
You can use the \mono{trace} and cursor keys to check the lower and upper class limits for each class as well as the count or frequency. The screenshot above indicates that the class of heights between 62 and 64 inches contains 7 individuals.%
\item\hypertarget{li-33}{}\hypertarget{p-60}{}%
If you use the \mono{trace} feature to determine the frequency for each bar (class) you can generate a \terminology{frequency distribution}. On the TI-84 Plus calculator, instead of using the \terminology{classes} we instead enter the \terminology{class midpoints} in list 2 and the \terminology{frequencies} in list 3.%
\begin{sidebyside}{2}{0}{0}{0}%
\begin{sbspanel}{0.5}%
{\centering%
\begin{tabular}{ll}\hrulethick
\multicolumn{1}{c}{Height}&\multicolumn{1}{c}{Freq.}\tabularnewline\hrulethick
\multicolumn{1}{c}{60\textendash{}\textless{}62}&\multicolumn{1}{c}{1}\tabularnewline\hrulethin
\multicolumn{1}{c}{62\textendash{}\textless{}64}&\multicolumn{1}{c}{7}\tabularnewline\hrulethin
\multicolumn{1}{c}{64\textendash{}\textless{}66}&\multicolumn{1}{c}{1}\tabularnewline\hrulethin
\multicolumn{1}{c}{66\textendash{}\textless{}68}&\multicolumn{1}{c}{2}\tabularnewline\hrulethin
\multicolumn{1}{c}{68\textendash{}\textless{}70}&\multicolumn{1}{c}{6}\tabularnewline\hrulethin
\multicolumn{1}{c}{70\textendash{}\textless{}72}&\multicolumn{1}{c}{5}\tabularnewline\hrulethin
\multicolumn{1}{c}{72\textendash{}\textless{}74}&\multicolumn{1}{c}{1}\tabularnewline\hrulethin
\multicolumn{1}{c}{74\textendash{}\textless{}76}&\multicolumn{1}{c}{0}\tabularnewline\hrulethin
\multicolumn{1}{c}{76\textendash{}\textless{}78}&\multicolumn{1}{c}{0}\tabularnewline\hrulethin
\multicolumn{1}{c}{78\textendash{}\textless{}80}&\multicolumn{1}{c}{1}\tabularnewline\hrulethick
\multicolumn{1}{r}{Sum:}&\multicolumn{1}{c}{24}\tabularnewline\hrulethick
\end{tabular}
\par}
\end{sbspanel}%
\begin{sbspanel}{0.5}%
\includegraphics[width=1\linewidth]{images/heights-freq-dist.png}
\end{sbspanel}%
\nopagebreak%
\begin{sbscaption}{0.5}%
\captionof{table}{Height frequency distribution\label{table-4}}
\end{sbscaption}%
\begin{sbscaption}{0.5}%
\end{sbscaption}%
\end{sidebyside}%
\item\hypertarget{li-34}{}\hypertarget{p-61}{}%
Next, we can have the TI-84 Plus calculator generate \terminology{relative frequencies} for us. Let's store the relative frequencies in list 4. Go to \mono{stat}, \mono{1} and edit list 4. Cursor to the list 4 header, i.e.\@ cursor up until the header is highlighted (shaded black) and press the following: \begin{sidebyside}{1}{0.2}{0.2}{0}%
\begin{sbspanel}{0.6}%
\hypertarget{p-62}{}%
\mono{2nd}\mono{3}\mono{÷}\mono{2nd}\mono{stat},%
\end{sbspanel}%
\end{sidebyside}%
 choose CALC and then press the following keys: \begin{sidebyside}{1}{0.2}{0.2}{0}%
\begin{sbspanel}{0.6}%
\hypertarget{p-63}{}%
\mono{5}\mono{2nd}\mono{3}\mono{)}\mono{enter}%
\end{sbspanel}%
\end{sidebyside}%
 \begin{sidebyside}{2}{0}{0}{0}%
\begin{sbspanel}{0.5}%
\includegraphics[width=1\linewidth]{images/generate-rel-freq.png}
\end{sbspanel}%
\begin{sbspanel}{0.5}%
\includegraphics[width=1\linewidth]{images/generated-rel-freq.png}
\end{sbspanel}%
\end{sidebyside}%
 The resulting frequency and relative frequency distiributions are reproduced here.%
\begin{table}
\centering
\begin{tabular}{lll}\hrulethick
\multicolumn{1}{c}{Height}&\multicolumn{1}{r}{Freq.}&\multicolumn{1}{r}{Rel. Freq.}\tabularnewline\hrulethick
\multicolumn{1}{c}{60\textendash{}\textless{}62}&\multicolumn{1}{r}{1}&\multicolumn{1}{r}{.04167}\tabularnewline\hrulethin
\multicolumn{1}{c}{62\textendash{}\textless{}64}&\multicolumn{1}{r}{7}&\multicolumn{1}{r}{.29167}\tabularnewline\hrulethin
\multicolumn{1}{c}{64\textendash{}\textless{}66}&\multicolumn{1}{r}{1}&\multicolumn{1}{r}{.04167}\tabularnewline\hrulethin
\multicolumn{1}{c}{66\textendash{}\textless{}68}&\multicolumn{1}{r}{2}&\multicolumn{1}{r}{.08333}\tabularnewline\hrulethin
\multicolumn{1}{c}{68\textendash{}\textless{}70}&\multicolumn{1}{r}{6}&\multicolumn{1}{r}{.25000}\tabularnewline\hrulethin
\multicolumn{1}{c}{70\textendash{}\textless{}72}&\multicolumn{1}{r}{5}&\multicolumn{1}{r}{.20833}\tabularnewline\hrulethin
\multicolumn{1}{c}{72\textendash{}\textless{}74}&\multicolumn{1}{r}{1}&\multicolumn{1}{r}{.04167}\tabularnewline\hrulethin
\multicolumn{1}{c}{74\textendash{}\textless{}76}&\multicolumn{1}{r}{0}&\multicolumn{1}{r}{.00000}\tabularnewline\hrulethin
\multicolumn{1}{c}{76\textendash{}\textless{}78}&\multicolumn{1}{r}{0}&\multicolumn{1}{r}{.00000}\tabularnewline\hrulethin
\multicolumn{1}{c}{78\textendash{}\textless{}80}&\multicolumn{1}{r}{1}&\multicolumn{1}{r}{.04167}\tabularnewline\hrulethick
\multicolumn{1}{r}{Sum:}&\multicolumn{1}{r}{24}&\multicolumn{1}{r}{1.00000}\tabularnewline\hrulethick
\end{tabular}
\caption{Height frequency\slash{}relative frequency distribution\label{table-5}}
\end{table}
\end{enumerate}
\end{subsectionptx}
%
%
\typeout{************************************************}
\typeout{Subsection 2.2.3 Creating Histograms from Frequency Distributions Using the TI-84 Plus}
\typeout{************************************************}
%
\begin{subsectionptx}{Creating Histograms from Frequency Distributions Using the TI-84 Plus}{}{Creating Histograms from Frequency Distributions Using the TI-84 Plus}{}{}{sub-histogram-from-freq-dist}
\hypertarget{p-64}{}%
Perhaps you don't have access to raw data, but you do have a frequency distribution of the data, then you can still create a histogram of the data.%
\par
\hypertarget{p-65}{}%
In this section, we will continue using the height data in table \hyperref[tab-student-heights]{Table~\ref{tab-student-heights}}. Further, we will assume you have the same list data generated in that example in lists 2,3 and 4 of your calculator.%
\par
\hypertarget{p-66}{}%
To create a histogram from a frequency distribution:%
\leavevmode%
\begin{enumerate}
\item\hypertarget{li-35}{}\hypertarget{p-67}{}%
Go to STAT PLOTS by pressing \mono{2nd} \mono{y=} and choosing Plot1. Assuming you have the class midpoints in list 2 and the frequencies in list 3, we can generate a histogram by filling out the Plot1 menu to look like below:%
\begin{figure}\centering\includegraphics[width=0.4\linewidth]{images/stat-plot-histogram-from-freq-dist.png}
\end{figure}\item\hypertarget{li-36}{}\hypertarget{p-68}{}%
If you have not changed any settings in \mono{window}, then you can just press \mono{graph} to generate the histogram. If you have changed window settings use ZoomStat by pressing \mono{zoom} \mono{9}.%
\begin{figure}\centering\includegraphics[width=0.4\linewidth]{images/histogram-from-freq-dist.png}
\end{figure}\end{enumerate}
\end{subsectionptx}
%
%
\typeout{************************************************}
\typeout{Reading Questions 2.2.4 Cigarette Tax Rates}
\typeout{************************************************}
%
\begin{reading-questions-subsection}{Cigarette Tax Rates}{}{Cigarette Tax Rates}{}{}{reading-questions-15}
\hypertarget{p-69}{}%
The table shows the tax in dollars, on a pack of cigarettes in each of the 50 states and Washington D.C. as of January 2014.%
\begin{table}
\centering
\begin{tabular}{lllllll}\hrulethick
0.425&1.339&0.60&0.17&0.45&1.53&2.52\tabularnewline\hrulethin
2.00&0.37&0.36&1.70&0.44&0.62&0.60\tabularnewline\hrulethin
2.00&3.20&2.00&0.64&1.25&1.41&2.50\tabularnewline\hrulethin
1.15&0.57&2.00&0.80&1.03&1.70&\tabularnewline\hrulethin
0.87&1.98&3.51&1.68&1.31&2.62&\tabularnewline\hrulethin
0.84&0.995&2.00&2.70&1.60&0.30&\tabularnewline\hrulethin
3.40&1.36&2.83&1.66&3.50&3.025&\tabularnewline\hrulethin
1.60&0.79&0.68&4.35&0.57&0.55&\tabularnewline\hrulethick
\end{tabular}
\caption{Cigarette tax rates in dollars (\textdollar{})\label{table-6}}
\end{table}
\hypertarget{p-70}{}%
Note: Normally when doing the following steps by hand, the first two steps are done in reverse, i.e.\@ you first create a \terminology{frequency distribution}, and then use that data to create a \terminology{histogram}. However, since we will use the TI graphing calculator for the first two steps, we will have the calculator do the counting for us and then just read the \terminology{frequency distribution} data from the histogram.%
\begin{divisionexercise}{1}{}{}{exercise-64}%
\hypertarget{p-71}{}%
Use your TI-84 Plus calculator to construct a histogram with lower limit \textdollar{}0.00, upper limit \textdollar{}5.00 and class width of \textdollar{}0.50.%
\textbf{Answer}.\quad%
\begin{figure}\centering\includegraphics[width=0.4\linewidth]{images/cigarette-tax-rates-histogram.png}
\end{figure}\end{divisionexercise}%
\begin{divisionexercise}{2}{}{}{exercise-65}%
\hypertarget{p-72}{}%
What is the shape of your distribution?%
\par\smallskip%
\noindent\textbf{Answer}.\hypertarget{answer-61}{}\quad%
skewed right\end{divisionexercise}%
\begin{divisionexercise}{3}{}{}{exercise-66}%
\hypertarget{p-73}{}%
Construct a frequency distribution and relative frequency distribution.%
\textbf{Hint}.\quad%
\hypertarget{p-74}{}%
Use your histogram to get the frequencies, put them in a list, and use your calculator to compute the relative frequencies.%
\textbf{Answer}.\quad%
\leavevmode%
\begin{sidebyside}{2}{0}{0}{0}%
\begin{sbspanel}{0.5}%
\includegraphics[width=1\linewidth]{images/cigarette-tax-rates-freq-dist.png}
\end{sbspanel}%
\begin{sbspanel}{0.5}%
{\centering%
\begin{tabular}{lll}\hrulethick
\multicolumn{1}{c}{Class}&\multicolumn{1}{r}{Freq.}&\multicolumn{1}{r}{Rel. Freq.}\tabularnewline\hrulethick
0.00\textemdash{}0.49&7&.13725\tabularnewline\hrulethin
0.50\textemdash{}0.99&13&.2549\tabularnewline\hrulethin
1.00\textemdash{}1.49&7&.13725\tabularnewline\hrulethin
1.50\textemdash{}1.99&8&.15686\tabularnewline\hrulethin
2.00\textemdash{}2.49&5&.09804\tabularnewline\hrulethin
2.50\textemdash{}2.99&5&.09804\tabularnewline\hrulethin
3.00\textemdash{}3.49&3&.05882\tabularnewline\hrulethin
3.50\textemdash{}3.99&2&.03922\tabularnewline\hrulethin
4.00\textemdash{}4.49&1&.01961\tabularnewline\hrulethin
4.50\textemdash{}4.99&0&.00000\tabularnewline\hrulethick
\end{tabular}
\par}
\end{sbspanel}%
\end{sidebyside}%
\end{divisionexercise}%
\begin{divisionexercise}{4}{}{}{exercise-67}%
\hypertarget{p-75}{}%
Construct a relative frequency histogram with \mono{window} settings:\leavevmode%
\begin{itemize}[label=\textbullet]
\item{}Xmin = 0%
\item{}Xmax = 5%
\item{}Xscl = 0.5%
\item{}Ymin = 0%
\item{}Ymax = 0.3%
\item{}Yscl = 0.05%
\end{itemize}
%
\par\smallskip%
\noindent\textbf{Answer}.\hypertarget{answer-63}{}\quad%
\begin{figure}\centering\includegraphics[width=0.4\linewidth]{images/cigarette-tax-rel-freq-histogram.png}
\end{figure}\end{divisionexercise}%
\begin{divisionexercise}{5}{}{}{exercise-68}%
\hypertarget{p-76}{}%
What percentage of states have a tax of \textdollar{}3 or more on a pack of cigarettes?%
\par\smallskip%
\noindent\textbf{Answer}.\hypertarget{answer-64}{}\quad%
\hypertarget{p-77}{}%
We can use the calculator to answer this:%
\begin{figure}\centering\includegraphics[width=0.4\linewidth]{images/cigarette-tax-rate-percentage.png}
\end{figure}\hypertarget{p-78}{}%
Thus about 11.8\% of states have a tax of \textdollar{}3 or more per pack of cigarettes.%
\end{divisionexercise}%
\end{reading-questions-subsection}
\begin{figure}
\centering
\setlength{\qrsize}{9em}
\setlength{\previewwidth}{\linewidth}
\addtolength{\previewwidth}{-\qrsize}
\begin{tcbraster}[raster columns=2, raster column skip=1pt, raster halign=center, raster force size=false, raster left skip=0pt, raster right skip=0pt]%
\begin{tcolorbox}[previewstyle, width=\previewwidth]%
\includegraphics[width=0.80\linewidth,height=\qrsize,keepaspectratio]{images/video-6.jpg}%
\end{tcolorbox}%
\begin{tcolorbox}[qrstyle]%
{\hypersetup{urlcolor=black}\qrcode[height=\qrsize]{https://www.youtube.com/watch?v=HMkllhBI91Y\&start=0\&end=218}}%
\end{tcolorbox}%
\begin{tcolorbox}[captionstyle]%
\small YouTube: \mono{https://www.youtube.com/watch?v=HMkllhBI91Y}\end{tcolorbox}%
\end{tcbraster}%
\caption{Crash Course Statistics: Visualization Part II\label{figure-7}}
\end{figure}
\begin{definition}{}{definition-51}%
\hypertarget{p-79}{}%
A \terminology{stem and leaf plot} is another way to display quantitative data where the rightmost or last digit of a datum is called a \terminology{leaf} and the other digits form the \terminology{stem}.%
\end{definition}
%
%
\typeout{************************************************}
\typeout{Reading Questions 2.2.5 Grams of Fat in McDonalds' Breakfast Menu Items}
\typeout{************************************************}
%
\begin{reading-questions-subsection}{Grams of Fat in McDonalds' Breakfast Menu Items}{}{Grams of Fat in McDonalds' Breakfast Menu Items}{}{}{reading-questions-16}
\hypertarget{p-80}{}%
%
\end{reading-questions-subsection}
\end{sectionptx}
%
%
\typeout{************************************************}
\typeout{Section 2.3 Graphical Misrepresentations of Data}
\typeout{************************************************}
%
\begin{sectionptx}{Graphical Misrepresentations of Data}{}{Graphical Misrepresentations of Data}{}{}{sec-graphical-misrepresentations-of-data}
\hypertarget{p-81}{}%
Placeholder text%
\end{sectionptx}
\end{chapterptx}
 %
%
\typeout{************************************************}
\typeout{Chapter 3 Numerically Summarizing Data}
\typeout{************************************************}
%
\begin{chapterptx}{Numerically Summarizing Data}{}{Numerically Summarizing Data}{}{}{ch-numerical-summaries}
%
%
\typeout{************************************************}
\typeout{Section 3.1 Measures of Central Tendency}
\typeout{************************************************}
%
\begin{sectionptx}{Measures of Central Tendency}{}{Measures of Central Tendency}{}{}{sec-measures-central-tendency}
\begin{definition}{Mean (Arithmetic Average).}{definition-52}%
\hypertarget{p-82}{}%
The \terminology{mean} or \terminology{arithmetic average} is computed by adding up all the values and dividing by the number of observations.%
\leavevmode%
\begin{description}
\item[{population mean}]\hypertarget{li-43}{} \(\mu = \dfrac{\sum x_i}{N}\) (pronounced ``mew'')\item[{sample mean}]\hypertarget{li-44}{}\hypertarget{p-83}{}%
\(\bar{x} = \dfrac{\sum x_i}{n}\) (pronounced ``x bar'')%
\end{description}
\hypertarget{p-84}{}%
Note that it is a convention in Statistics to use Greek letters to represent population statistics while we use Roman letters to denote sample statistics.%
\leavevmode%
\begin{itemize}[label=\textbullet]
\item{}\alert{Advantage:} the mean is used as a starting point for computing several other useful statistics e.g.\@ std. deviation.%
\item{}\alert{Disadvantage:} If the data is skewed or contains outliers, then the mean is not the best way to measure the center of your data.%
\end{itemize}
\end{definition}
\begin{definition}{Median.}{definition-53}%
\hypertarget{p-85}{}%
The \terminology{median} or \terminology{middle value} is the value that lies in the middle of the data when the data is arranged in ascending order. Notation: \(M\).%
\par
\hypertarget{p-86}{}%
Algorithm:\leavevmode%
\begin{enumerate}
\item\hypertarget{li-47}{}Sort the data, from least to greatest.%
\item\hypertarget{li-48}{}Find the single middle element, or middle two elements if the list has an even number of elements.%
\item\hypertarget{li-49}{}If there is one value in the middle (\(n\) odd), then that value is the median. If there are two values in the middle (\(n\) even) then the median is the average of the two middle values.%
\end{enumerate}
%
\par
\hypertarget{p-87}{}%
\leavevmode%
\begin{itemize}[label=\textbullet]
\item{}\alert{Advantage:} Resistant to outliers and or skewed data.%
\item{}\alert{Disadvantage:} Hard to compute by hand when you have a large data set.%
\end{itemize}
%
\end{definition}
\begin{example}{Computing the Median By Hand I.}{example-1}%
\hypertarget{p-88}{}%
Find the median of the list: 3, 4, 2, 6, 3.%
\par\smallskip%
\noindent\textbf{Solution}.\hypertarget{solution-8}{}\quad%
\hypertarget{p-89}{}%
The ordered list is: 2, 3, 3, 4, 6, and the middle value is 3, thus \(M=3\).%
\end{example}
\begin{example}{Computing the Median By Hand II.}{example-2}%
\hypertarget{p-90}{}%
Find the median of the list: 3, 4, 2, 6, 3, 7%
\par\smallskip%
\noindent\textbf{Solution}.\hypertarget{solution-9}{}\quad%
\hypertarget{p-91}{}%
The ordered list is: 2, 3, 3, 4, 6, 7, but it has an even number of elements, so we compute the average or mean of the middle two elements, thus the median is:%
\begin{gather*}
M = \frac{3+4}{2} = 3.5
\end{gather*}
%
\end{example}
\begin{definition}{Resistant.}{definition-54}%
\hypertarget{p-92}{}%
A statistic is said to be \terminology{resistant} if extreme values do not affect the value of the statistic.%
\end{definition}
\begin{definition}{Mode.}{definition-55}%
\hypertarget{p-93}{}%
The \terminology{mode} is the most frequently occuring observation in a data set. A single data set may have multiple modes.%
\end{definition}
\begin{figure}
\centering
\setlength{\qrsize}{9em}
\setlength{\previewwidth}{\linewidth}
\addtolength{\previewwidth}{-\qrsize}
\begin{tcbraster}[raster columns=2, raster column skip=1pt, raster halign=center, raster force size=false, raster left skip=0pt, raster right skip=0pt]%
\begin{tcolorbox}[previewstyle, width=\previewwidth]%
\includegraphics[width=0.80\linewidth,height=\qrsize,keepaspectratio]{images/video-7.jpg}%
\end{tcolorbox}%
\begin{tcolorbox}[qrstyle]%
{\hypersetup{urlcolor=black}\qrcode[height=\qrsize]{https://www.youtube.com/watch?v=kn83BA7cRNM}}%
\end{tcolorbox}%
\begin{tcolorbox}[captionstyle]%
\small YouTube: \mono{https://www.youtube.com/watch?v=kn83BA7cRNM}\end{tcolorbox}%
\end{tcbraster}%
\caption{Crash Course Statistics: Mean, Median and Mode: Measures of Central Tendency\label{figure-8}}
\end{figure}
%
%
\typeout{************************************************}
\typeout{Reading Questions 3.1 Working Hours}
\typeout{************************************************}
%
\begin{reading-questions-subsection-numberless}{Working Hours}{}{Working Hours}{}{}{reading-questions-17}
\hypertarget{p-94}{}%
A random sample of 25 college students was asked, ``How many hours per week do you typically work outside your home?'' Their responses are recorded in the table below.%
\begin{table}
\centering
\begin{tabular}{lllll}\hrulethick
\multicolumn{1}{r}{0}&\multicolumn{1}{r}{32}&\multicolumn{1}{r}{15}&\multicolumn{1}{r}{20}&\multicolumn{1}{r}{30}\tabularnewline\hrulethin
\multicolumn{1}{r}{40}&\multicolumn{1}{r}{30}&\multicolumn{1}{r}{20}&\multicolumn{1}{r}{35}&\multicolumn{1}{r}{35}\tabularnewline\hrulethin
\multicolumn{1}{r}{28}&\multicolumn{1}{r}{38}&\multicolumn{1}{r}{36}&\multicolumn{1}{r}{25}&\multicolumn{1}{r}{25}\tabularnewline\hrulethin
\multicolumn{1}{r}{30}&\multicolumn{1}{r}{5}&\multicolumn{1}{r}{28}&\multicolumn{1}{r}{30}&\multicolumn{1}{r}{24}\tabularnewline\hrulethin
\multicolumn{1}{r}{28}&\multicolumn{1}{r}{30}&\multicolumn{1}{r}{35}&\multicolumn{1}{r}{15}&\multicolumn{1}{r}{32}\tabularnewline\hrulethick
\end{tabular}
\caption{Working hours of 25 college students\label{table-7}}
\end{table}
\begin{divisionexercise}{1}{}{}{exercise-69}%
\hypertarget{p-95}{}%
Determine the distribution of the hours worked by creating a frequency histogram on your TI-84 Plus calculator. Use 0 as your lower class limit and 5 as your class width.%
\par\smallskip%
\noindent\textbf{Answer}.\hypertarget{answer-65}{}\quad%
\hypertarget{p-96}{}%
First enter the data into any list, press \mono{window} and enter the following values:\leavevmode%
\begin{itemize}[label=\textbullet]
\item{}Xmin=0%
\item{}Xmax=45%
\item{}Xscl=5%
\item{}Ymin=0%
\item{}Ymax=10%
\item{}Yscl=2%
\end{itemize}
Finally after choosing the correct list in STAT PLOT press \mono{graph} to plot the histogram with the above settings.%
\begin{figure}\centering\includegraphics[width=0.5\linewidth]{images/hours-worked-freq-histogram.png}
\end{figure}\end{divisionexercise}%
\begin{divisionexercise}{2}{}{}{exercise-70}%
\hypertarget{p-97}{}%
Use your graphing calculator to find the mean and median. Which measure better describes the hours worked?%
\par\smallskip%
\noindent\textbf{Answer}.\hypertarget{answer-66}{}\quad%
\leavevmode%
\begin{sidebyside}{2}{0}{0}{0}%
\begin{sbspanel}{0.5}%
\includegraphics[width=1\linewidth]{images/hours-worked-mean.png}
\end{sbspanel}%
\begin{sbspanel}{0.5}%
\includegraphics[width=1\linewidth]{images/hours-worked-median.png}
\end{sbspanel}%
\end{sidebyside}%
\par
\hypertarget{p-98}{}%
\leavevmode%
\begin{itemize}[label=\textbullet]
\item{}\(\bar{x} = 26.64\)%
\item{}\(M = 30\)%
\end{itemize}
Because the data is skewed to the left, the median is a better measure of the central tendency of the data.%
\end{divisionexercise}%
\end{reading-questions-subsection-numberless}
\end{sectionptx}
%
%
\typeout{************************************************}
\typeout{Section 3.2 Measures of Dispersion}
\typeout{************************************************}
%
\begin{sectionptx}{Measures of Dispersion}{}{Measures of Dispersion}{}{}{sec-measures-dispersion}
\hypertarget{p-99}{}%
The video below mentions the inter-quartile range or IQR which we cover in section \hyperref[sec-measures-position-and-outliers]{Section~\ref{sec-measures-position-and-outliers}}. The IQR is simply the middle 50\% of the data.%
\begin{figure}
\centering
\setlength{\qrsize}{9em}
\setlength{\previewwidth}{\linewidth}
\addtolength{\previewwidth}{-\qrsize}
\begin{tcbraster}[raster columns=2, raster column skip=1pt, raster halign=center, raster force size=false, raster left skip=0pt, raster right skip=0pt]%
\begin{tcolorbox}[previewstyle, width=\previewwidth]%
\includegraphics[width=0.80\linewidth,height=\qrsize,keepaspectratio]{images/video-8.jpg}%
\end{tcolorbox}%
\begin{tcolorbox}[qrstyle]%
{\hypersetup{urlcolor=black}\qrcode[height=\qrsize]{https://www.youtube.com/watch?v=R4yfNi_8Kqw}}%
\end{tcolorbox}%
\begin{tcolorbox}[captionstyle]%
\small YouTube: \mono{https://www.youtube.com/watch?v=R4yfNi\_8Kqw}\end{tcolorbox}%
\end{tcbraster}%
\caption{Crash Course Statistics: Measures of Spread\label{figure-9}}
\end{figure}
\end{sectionptx}
%
%
\typeout{************************************************}
\typeout{Section 3.3 Measures of Central Tendency and Dispersion from Grouped Data}
\typeout{************************************************}
%
\begin{sectionptx}{Measures of Central Tendency and Dispersion from Grouped Data}{}{Measures of Central Tendency and Dispersion from Grouped Data}{}{}{sec-measures-grouped-data}
\hypertarget{p-100}{}%
Placeholder text%
\end{sectionptx}
%
%
\typeout{************************************************}
\typeout{Section 3.4 Measures of Position and Outliers}
\typeout{************************************************}
%
\begin{sectionptx}{Measures of Position and Outliers}{}{Measures of Position and Outliers}{}{}{sec-measures-position-and-outliers}
\begin{definition}{Z-score.}{definition-56}%
\hypertarget{p-101}{}%
A \terminology{z-score} of a data point is the number of standard deviations the point is from the mean.%
\leavevmode%
\begin{description}
\item[{population z-score}]\hypertarget{li-60}{}\hypertarget{p-102}{}%
\(z = \dfrac{x-\mu}{\sigma}\)%
\item[{sample z-score}]\hypertarget{li-61}{}\hypertarget{p-103}{}%
\(z = \dfrac{x-\bar{x}}{s}\)%
\end{description}
 Characteristics of z-scores: \leavevmode%
\begin{enumerate}
\item\hypertarget{li-62}{}unitless%
\item\hypertarget{li-63}{}mean of z-scores = 0%
\item\hypertarget{li-64}{}std. deviation of z-scores=1%
\item\hypertarget{li-65}{}If a datum is larger than the mean its z-score is positive.%
\item\hypertarget{li-66}{}If a datum is smaller than the mean its z-score is negative.%
\item\hypertarget{li-67}{}Z-scores can be used to compare measurements on different scales.%
\end{enumerate}
\end{definition}
\begin{figure}
\centering
\setlength{\qrsize}{9em}
\setlength{\previewwidth}{\linewidth}
\addtolength{\previewwidth}{-\qrsize}
\begin{tcbraster}[raster columns=2, raster column skip=1pt, raster halign=center, raster force size=false, raster left skip=0pt, raster right skip=0pt]%
\begin{tcolorbox}[previewstyle, width=\previewwidth]%
\includegraphics[width=0.80\linewidth,height=\qrsize,keepaspectratio]{images/video-9.jpg}%
\end{tcolorbox}%
\begin{tcolorbox}[qrstyle]%
{\hypersetup{urlcolor=black}\qrcode[height=\qrsize]{https://www.youtube.com/watch?v=uAxyI_XfqXk}}%
\end{tcolorbox}%
\begin{tcolorbox}[captionstyle]%
\small YouTube: \mono{https://www.youtube.com/watch?v=uAxyI\_XfqXk}\end{tcolorbox}%
\end{tcbraster}%
\caption{Crash Course Statistics: Z-Scores and Percentiles\label{figure-10}}
\end{figure}
\begin{example}{Comparing Final Exam Scores.}{example-3}%
\hypertarget{p-104}{}%
Final exam scores for a Biology course have a mean of 79 and a standard deviation of 5. Final exam scores for a Statistics course have a mean of 82 and a standard deviation of 7. You scored 86 on the Biology exam and 91 on the Statistics exam. How many standard deviations above the mean were you in each course? Relative to the rest of the students, on which exam did you do better?%
\par\smallskip%
\noindent\textbf{Solution}.\hypertarget{solution-10}{}\quad%
\begin{tabular}{llll}\hrulethick
\multicolumn{1}{r}{statistic}&\multicolumn{1}{c}{symbol}&\multicolumn{1}{c}{Biology}&\multicolumn{1}{c}{Statistics}\tabularnewline\hrulethick
\multicolumn{1}{r}{test score}&\multicolumn{1}{c}{\(x\)}&\multicolumn{1}{c}{86}&\multicolumn{1}{c}{91}\tabularnewline\hrulethin
\multicolumn{1}{r}{mean}&\multicolumn{1}{c}{\(\mu\)}&\multicolumn{1}{c}{79}&\multicolumn{1}{c}{82}\tabularnewline\hrulethin
\multicolumn{1}{r}{std. dev.}&\multicolumn{1}{c}{\(\sigma\)}&\multicolumn{1}{c}{5}&\multicolumn{1}{c}{7}\tabularnewline\hrulethin
\multicolumn{1}{r}{z-score}&\multicolumn{1}{c}{\(z\)}&\multicolumn{1}{c}{\(\dfrac{86-79}{5}\)}&\multicolumn{1}{c}{\(\dfrac{91-82}{7}\)}\tabularnewline\hrulethick
&\multicolumn{1}{c}{\(z\)}&\multicolumn{1}{c}{1.4}&\multicolumn{1}{c}{\(\approx\) 1.3}\tabularnewline\hrulethick
\end{tabular}
\hypertarget{p-105}{}%
Thus since, 1.4 \textgreater{} 1.3, you did slightly better on the Biology exam than you did on your Statistics exam, relative to your fellow students. This is despite the fact that in absolute terms, your score in Statistics was higher.%
\par
\hypertarget{p-106}{}%
This example demonstrates that it is meaningless to compare absolute measures which come from different distributions, in this case, the test scores. However, by rescaling each test score relative to the mean and the standard deviation of each distribution i.e.\@ creating z-scores, we can meaningfully compare data from different distributions.%
\end{example}
\begin{example}{Triathlon Times.}{example-4}%
\hypertarget{p-107}{}%
Roberto finishes a triathlon (750 m swim, 5 km run, and 20 km bicycle) in 63.2 minutes. Among all men in the race, the mean finishing time was 69.4 minutes with a standard deviation of 8.0 minutes. Zandra finishes the same triathlon in 79.3 minutes. Among all women in the race, the mean finishing time was 84.7 minutes with a standard deviation of 7.4 minutes. Who did better relative to their gender?%
\par\smallskip%
\noindent\textbf{Solution}.\hypertarget{solution-11}{}\quad%
\hypertarget{p-108}{}%
We compute a z-score for each athlete:%
\begin{tabular}{lll}\hrulethick
\multicolumn{1}{r}{statistic}&\multicolumn{1}{c}{Roberto}&\multicolumn{1}{c}{Zandra}\tabularnewline\hrulethin
\multicolumn{1}{r}{z-score}&\multicolumn{1}{c}{\(\dfrac{63.2-69.4}{8.0}\)}&\multicolumn{1}{c}{\(\dfrac{79.3-84.7}{7.4}\)}\tabularnewline\hrulethin
\multicolumn{1}{r}{z-score}&\multicolumn{1}{c}{-0.775}&\multicolumn{1}{c}{\(\approx\) -0.730}\tabularnewline\hrulethick
\end{tabular}
\hypertarget{p-109}{}%
In the case of a race, it is better to have a time below the mean time of all the racers, thus the more negative the z-score the better. Since Roberto has the lower z-score, Roberto did better relative to his competitors than Zandra did relative to hers.%
\end{example}
\end{sectionptx}
%
%
\typeout{************************************************}
\typeout{Section 3.5 The Five-{}-{}Number Summary and Boxplots}
\typeout{************************************************}
%
\begin{sectionptx}{The Five-{}-{}Number Summary and Boxplots}{}{The Five-{}-{}Number Summary and Boxplots}{}{}{sec-five-number-summary-boxplots}
\hypertarget{p-110}{}%
Placeholder text%
\end{sectionptx}
\end{chapterptx}
 %
%
\typeout{************************************************}
\typeout{Chapter 4 Describing the Relation between Two Variables}
\typeout{************************************************}
%
\begin{chapterptx}{Describing the Relation between Two Variables}{}{Describing the Relation between Two Variables}{}{}{ch-relations-between-two-varibles}
\end{chapterptx}
\end{partptx}
%
%
\typeout{************************************************}
\typeout{Part III Probability and Probability Distributions}
\typeout{************************************************}
%
\begin{partptx}{Probability and Probability Distributions}{}{Probability and Probability Distributions}{}{}{part-probability}
 %
%
\typeout{************************************************}
\typeout{Chapter 5 Probability}
\typeout{************************************************}
%
\begin{chapterptx}{Probability}{}{Probability}{}{}{ch-probability}
\end{chapterptx}
 %
%
\typeout{************************************************}
\typeout{Chapter 6 Discrete Probability Distributions}
\typeout{************************************************}
%
\begin{chapterptx}{Discrete Probability Distributions}{}{Discrete Probability Distributions}{}{}{ch-discrete-probability}
\end{chapterptx}
 %
%
\typeout{************************************************}
\typeout{Chapter 7 The Normal Probability Distribution}
\typeout{************************************************}
%
\begin{chapterptx}{The Normal Probability Distribution}{}{The Normal Probability Distribution}{}{}{ch-normal-probability}
\end{chapterptx}
\end{partptx}
%
%
\typeout{************************************************}
\typeout{Part IV Inference: From Samples to Population}
\typeout{************************************************}
%
\begin{partptx}{Inference: From Samples to Population}{}{Inference: From Samples to Population}{}{}{part-inferential-stats}
 %
%
\typeout{************************************************}
\typeout{Chapter 8 Sampling Distributions}
\typeout{************************************************}
%
\begin{chapterptx}{Sampling Distributions}{}{Sampling Distributions}{}{}{ch-sampling-distributions}
\end{chapterptx}
 %
%
\typeout{************************************************}
\typeout{Chapter 9 Estimating the Value of a Parameter}
\typeout{************************************************}
%
\begin{chapterptx}{Estimating the Value of a Parameter}{}{Estimating the Value of a Parameter}{}{}{ch-estimating-parameters}
\end{chapterptx}
 %
%
\typeout{************************************************}
\typeout{Chapter 10 Hypothesis Tests Regarding a Parameter}
\typeout{************************************************}
%
\begin{chapterptx}{Hypothesis Tests Regarding a Parameter}{}{Hypothesis Tests Regarding a Parameter}{}{}{ch-hypothesis-tests}
\end{chapterptx}
 %
%
\typeout{************************************************}
\typeout{Chapter 11 Inferences on Two Samples}
\typeout{************************************************}
%
\begin{chapterptx}{Inferences on Two Samples}{}{Inferences on Two Samples}{}{}{ch-inferences-on-two-samples}
\end{chapterptx}
\end{partptx}
\end{document}